\subsection{用数值密度定义的测度}

定义 2. --- 设 f 是定义在 T 上、取值于 \(\overline{R}\) 或 Banach 空间的函数。若称 f 是局部 \(\mu\) -可积的，意指 f 是 \(\mu\) -可测的，并且 T 中每一点 \(x \in T\) 都有一个邻域 V，使得 \(\mu^{\bullet}(|\mathbf{f}|\varphi_{\mathrm{V}}) < +\infty\) 。

这个定义在 T 为局部紧的情况下与 Ch. V, §5, No.~1 中给出的定义一致。

设 \(f\) 是一个局部 \(\mu\) -可积的正函数；映射 \(K \mapsto f_K \cdot \mu_K\) 是一个预测度（Ch. V, §7, No.~1, Cor. 2 of Th. 1），将其记作 \(f \cdot \mu\) 。

命题 2. --- 若 \(f\) 是正的局部 \(\mu\) -可积函数，则对于每个函数 \(g \in \mathcal{F}_{+}(\mathrm{T})\) ，都有关系式

\[
(f \cdot \mu) ^ {\bullet} (g) = \mu^ {\bullet} (f g).\tag{1}
\]

事实上，对 T 中的每个紧集 K，

\[
(f \cdot \mu) _ {\mathrm{K}} ^ {\bullet} (g _ {\mathrm{K}}) = (f _ {\mathrm{K}} \cdot \mu_ {\mathrm{K}}) ^ {\bullet} (g _ {\mathrm{K}}) = \mu_ {\mathrm{K}} ^ {\bullet} (f _ {\mathrm{K}} g _ {\mathrm{K}}) = \mu_ {\mathrm{K}} ^ {\bullet} ((f g) _ {\mathrm{K}}),
\]

利用 \(f \cdot \mu\) 的定义以及 Ch. V, §5, No.~3 的 Prop. 3。将 \(K\) 取遍并取上确界，即得 Prop. 2。

现在令 \(x \in \mathbf{T}\)，并令 \(\mathbf{V}\) 是 \(x\) 的一个邻域，使得 \(\mu^{\bullet}(f\varphi_{\mathrm{V}}) < +\infty\)（Def. 2）；于是 \((f \cdot \mu)^{\bullet}(\mathrm{V}) = \mu^{\bullet}(f\varphi_{\mathrm{V}}) < +\infty\) ，因此 \(f \cdot \mu\) 是一个测度。

定义 3. --- 设 \(f\) 是一个局部 \(\mu\) -可积的正函数。测度 \(f \cdot \mu : K \mapsto f_K \cdot \mu_K\) 称为具有密度 \(f\) 且相对于 \(\mu\) 的测度，或 \(\mu\) 与函数 f 的乘积测度。每个形如 \(f \cdot \mu\) 的测度（其中 f 为正的局部 \(\mu\) -可积函数）称为以 \(\mu\) 为基的测度。

注记。--- 1) \(f \cdot \mu\) 的定义可以推广到 \(f\) 是复的局部可积函数的情形；此时有 \(|f \cdot \mu| = |f| \cdot \mu\) ，这立即蕴含 \(f \cdot \mu\) 是测度而不只是预测度。我们保留“以 \(\mu'\) 为基的测度”这一表述，用来指称如此定义的复测度。

\begin{enumerate}
\def\labelenumi{\arabic{enumi})}
\setcounter{enumi}{1}
\tightlist
\item
  类似地，若 \(\theta\) 是一个复测度，则称 \(f\) 局部 \(\theta\) -可积，如果 f 局部 \(|\theta|\) -可积；并定义测度 \(f \cdot \theta : K \mapsto f_K \cdot \theta_K\) 。有 \(|f \cdot \theta| = |f| \cdot |\theta|\)（Ch. V, §5, No.~2, Prop. 2）。本号中，我们将撇开一切关于非正测度的内容。
\end{enumerate}

命题 3. --- 设 \(\nu\) 是 T 上的测度。要使 \(\nu\) 形如 \(f \cdot \mu\)，其中 \(f\) 是局部 \(\mu\) -可积的正函数，必要且充分的条件是每个 \(\mu\) -可忽略的紧集都是 \(\nu\) -可忽略的。若 \(f'\) 是另一个局部 \(\mu\) -可积的函数且 \(\nu = f' \cdot \mu\)，则 \(f = f'\) 局部 \(\mu\) -几乎处处成立。

这个条件显然是必要的（Prop. 2）。反之，假设每个 \(\mu\) -可忽略的紧集都是 \(\nu\) -可忽略的。引入 \((\mathrm{K}_{\alpha})_{\alpha \in \mathrm{A}}\) 作为 \(\mathbf{T}\) 关于测度 \(\mu + \nu\) 的一个分割，并令 \(\mathbf{N} = \mathbf{T} - \bigcup_{\alpha \in \mathrm{A}} \mathrm{K}_{\alpha}\) 。显然，\((\mathrm{K}_{\alpha})_{\alpha \in \mathrm{A}}\) 也是 \(\mu\) 和 \(\nu\) 的分割，因此 §1, No.~8 的 Prop. 9 蕴含，对于每个 \(g \in \mathcal{F}_+\) 有以下关系：

\[
\mu^ {\bullet} (g) = \sum_ {\alpha \in A} \mu_ {K _ {\alpha}} ^ {\bullet} \left(g _ {K _ {\alpha}}\right), \quad \nu^ {\bullet} (g) = \sum_ {\alpha \in A} \nu_ {K _ {\alpha}} ^ {\bullet} \left(g _ {K _ {\alpha}}\right).
\]

考虑一个紧集 \(C \subset K_{\alpha}\)，它对于 \(\mu_{K_{\alpha}}\) 是可忽略的；则 C 在局部意义下对于 \(\mu\) 可忽略，从而对于 \(\nu\) 可忽略，最后按照 \(\nu_{K_{\alpha}}\) 的定义，它对于 \(\nu\) 可忽略。由 Lebesgue--Nikodym 定理（Ch. V, §5, No.~5, Th. 2）可知，\(\nu_{K_{\alpha}}\) 关于 \(f_{\alpha}\) 具有一个密度 \(\mu_{K_{\alpha}}\)。令 f 为如下函数：在每个集合 \(f_{\alpha}\) 上与 \(K_{\alpha}\) 一致，在 N 上为 0；则 f 是 \(\mu\) -可测的（Ch. IV, §5, No.~10, Prop. 16），并且对于每个函数 \(g \in F_{+}\)，由上述关系及 Ch. V, §5, No.~3 的 Prop. 3，有

\[
\nu^ {\bullet} (g) = \sum_ {\alpha \in \mathrm{A}} \nu_ {\mathrm{K} _ {\alpha}} ^ {\bullet} \left(g _ {\mathrm{K} _ {\alpha}}\right) = \sum_ {\alpha \in \mathrm{A}} \mu_ {\mathrm{K} _ {\alpha}} ^ {\bullet} \left(f _ {\alpha} g _ {\mathrm{K} _ {\alpha}}\right) = \sum_ {\alpha \in \mathrm{A}} \mu_ {\mathrm{K} _ {\alpha}} ^ {\bullet} ((f g) _ {\mathrm{K} _ {\alpha}}) = \mu^ {\bullet} (f g).
\]

首先可得 \(f\) 局部 \(\mu\) -可积：若 \(x\) 是 \(T\) 中的一点，且 \(V\) 是 \(x\) 的一个邻域并满足 \(\nu^{\bullet}(V) < +\infty\)，则 \(\mu^{\bullet}(f\varphi_V) < +\infty\)。其次，Prop. 2 表明测度 \(\nu\) 与 \(f \cdot \mu\) 具有相同的本质上积分。因此二者相等（§1, No.~2, Cor. of Prop. 2）。\(f\) 的唯一性由紧空间的情形显然可得，命题得证。

注记 3). --- 关于测度 \(\nu = f \cdot \mu\) 的积分理论立即归结为 Ch. V 中处理的理论。为此，令 \((\mathrm{K}_{\alpha})_{\alpha \in \mathbf{A}}\) 是 T 关于 \(\mu\) 的一个分割，因而也是关于 \(\nu\) 的分割，并令 \(\mathrm{T}'\) 是 §1, No.~8 的 Scholium 中定义的局部紧空间；我们把 \(\mu\)（相应地，\(\nu\)）关联到测度 \(\mu'\)（相应地，\(\nu'\)），这些测度定义在 \(\mathrm{T}'\) 上，使得对于 \(\mu\) 和 \(\mu'\)（相应地，对于 \(\nu\) 和 \(\nu'\)），可测函数、取值于 Banach 空间的本质可积函数以及正函数的本质上积分都相同。因此 \(f\) 是 \(\mu'\) -可测的；f 是局部 \(\mu'\) -可积的，因为 \(\mathrm{T}'\) 是局部紧的，并且 \(\mathrm{T}'\) 的一个紧子集只与有限个紧集 \(\mathrm{K}_{\alpha}\) 相交（\(\alpha \in \mathbf{A}\)）。最后，关系 \(\nu'^{\bullet}(g) = \nu^{\bullet}(g) = \mu^{\bullet}(fg) = \mu'^{\bullet}(fg)\) 证明 \(\nu' = f \cdot \mu'\)（Ch. V, §5, No.~3, Prop. 3）。把 Ch. V, §5 的结果转写出来的工作留给读者。

\subsection{测度的像}

定义 4. --- 设 \(\pi\) 是从 T 到拓扑空间 X 的映射。称 \(\pi\) 为 \(\mu\) -逆紧的，是指 \(\pi\) 是 \(\mu\) -可测的，并且 X 中每一点 \(x\) 都有一个邻域 V，使得 \(\mu^{\bullet}\left(\overline{\pi}^{-1}(\mathrm{V})\right) < +\infty\) 。

注记。-- 1) 当 T 和 X 都是局部紧空间时，这个定义等价于 Ch. V, §6, No.~1 中的定义。

\begin{enumerate}
\def\labelenumi{\arabic{enumi})}
\setcounter{enumi}{1}
\item
  T 到 X 的适当连续映射（GT, I, §10, No.~1, Def. 1）对于每个测度 \(\mu\) 都是 \(\mu\) -逆紧的。事实上，令 \(x \in X\)；由于 \(\overline{\pi}^{-1}(x)\) 是紧的（loc. cit., No.~2, Th. 1），集合 \(\overline{\pi}^{-1}(x)\) 有一个开邻域 H，使得 \(\mu^{\bullet}(\mathrm{H}) < +\infty\)。令 \(\mathrm{V} = \mathrm{X} - \pi(\mathrm{T} - \mathrm{H})\)；由于 \(\pi\) 是闭映射，V 是 X 中的开集，包含 \(x\)，并满足 \(\overline{\pi}^{-1}(\mathrm{V}) \subset \mathrm{H}\)，从而 \(\mu^{\bullet}(\overline{\pi}^{-1}(\mathrm{V})) \leqslant \mu^{\bullet}(\mathrm{H}) < +\infty\)。
\item
  如果 \(\mu\) 有界，则从 T 到 X 的每个 \(\mu\) -可测映射都是 \(\mu\) -逆紧的。
\item
  如果 \(\theta\) 是 \(\mathbf{T}\) 上的复测度，则称 \(\pi\) 是 \(\theta\) -逆紧的，如果 \(\pi\) 对正测度 \(|\theta|\) 是逆紧的。
\end{enumerate}

命题 4. --- 设 \(\pi\) 是从 T 到拓扑空间 X 的 \(\mu\) -逆紧映射。X 上存在唯一的测度 \(\nu\)，使得 \(\nu^{\bullet}\) 等于像负荷 \(\pi(\mu^{\bullet})\)（§1, No.~1），换言之，使得 \(\nu^{\bullet}(g) = \mu^{\bullet}(g \circ \pi)\) 对所有 \(g \in \mathcal{F}_{+}(\mathbf{X})\) 成立。

唯一性显然（§1, No.~2, Cor. of Prop. 2）。为证明存在性，先处理 \(\mu\) 集中于紧集 K 且 \(\pi\) 在 K 上的限制连续的情形。此时 \(L = \pi(K)\) 是紧的；令 \(\pi'\) 为由 \(\pi\) 诱导的从 K 到 L 的连续映射，并令 \(\nu'\) 为 L 上的像测度 \(\pi'(\mu_K)\)，令 \(\nu\) 为由 \(\nu'\) 在 X 上定义的测度（§1, No.~3, Example 2）。对于每个 \(g \in \mathcal{F}_+(X)\)，

\[
\nu^ {\bullet} (g) = \nu^ {\prime \bullet} (g _ {\mathrm{L}}) = \mu_ {\mathrm{K}} ^ {\bullet} (g _ {\mathrm{L}} \circ \pi^ {\prime}) = \mu_ {\mathrm{K}} ^ {\bullet} ((g \circ \pi) _ {\mathrm{K}}) = \mu^ {\bullet} ((g \circ \pi) _ {\mathrm{K}} ^ {0}) = \mu^ {\bullet} (g \circ \pi)
\]

（我们依次使用了 §1, No.~3 的公式 (3)、Ch. V, §6, No.~2 的 Prop. 2、\(\mu_{\mathrm{K}}^{\bullet}\) 的定义以及 \(\mu\) 集中于 K 这一事实）。换言之，\(\nu^{\bullet} = \pi(\mu^{\bullet})\)。

现在转到一般情形；根据 §1, No.~8 的 Props. 10 和 9，\(\mu\) 是一个以紧支集测度 \((\mu_{\alpha})_{\alpha \in \mathbf{A}}\) 的可求和族之和，并且 \(\pi\) 在 \(\mathbf{K}_{\alpha}\)（测度 \(\mu_{\alpha}\) 的支集）上的限制对于每个 \(\alpha \in \mathbf{A}\) 都是连续的。上面处理的特殊情形允许对定义在 \(\mu_{\alpha}\) 上的每个测度 \(\mathbf{T}\) 关联定义在 \(\nu_{\alpha}\) 上的测度 \(\mathbf{X}\)，使得 \(\nu_{\alpha}^{\bullet} = \pi (\mu_{\alpha}^{\bullet})\)。于是，对于 \(g \in \mathcal{F}_{+}(\mathbf{X})\)，

\[
\sum_ {\alpha \in A} \nu_ {\alpha} ^ {\bullet} (g) = \sum_ {\alpha \in A} \mu_ {\alpha} ^ {\bullet} (g \circ \pi) = \mu^ {\bullet} (g \circ \pi).
\]

像负荷 \(\pi (\mu^{\bullet})\) 是局部有界的，因为 \(\pi\) 是 \(\mu\) -逆紧的；因此族 \((\nu_{\alpha})_{\alpha \in \mathbf{A}}\) 可求和（§1, No.~7, Prop. 7），其和 \(\nu\) 满足所述结论。

定义 5. --- 如果 \(\pi\) 是从 T 到拓扑空间 X 的 \(\mu\) -逆紧映射，则称 X 上的唯一测度 \(\nu\) 满足 \(\nu^{\bullet}(g) = \mu^{\bullet}(g \circ \pi)\) 对所有 \(g \in \mathcal{F}_{+}(\mathrm{X})\) 成立时，为 \(\mu\) 在 \(\pi\) 下的像测度，记作 \(\pi(\mu)\)。

例。--- 设 K 是 T 的紧子空间，i 是 K 到 T 的典范注入，\(\lambda\) 是 K 上的测度；由于 i 连续且 \(\lambda\) 有界，i 是 \(\lambda\) -逆紧的。§1, No.~3 的公式 (3) 表明“由 \(\lambda\) 在 T 上定义的测度”就是像测度 \(i(\lambda)\)。

注记 5). --- 如果 \(\theta\) 是实测度且 \(\pi\) 是 \(\theta\) -逆紧的，则 \(\pi\) 对测度 \(\theta^{+}\) 和 \(\theta^{-}\) 都是逆紧的；于是置 \(\pi(\theta) = \pi(\theta^{+}) - \pi(\theta^{-})\)。如果 \(\theta\) 是复测度且 \(\pi\) 是 \(\theta\) -逆紧的，则 \(\pi\) 对实测度 \(\mathcal{R}(\theta)\) 和 \(\mathcal{I}(\theta)\) 是逆紧的；于是置

\[
\pi (\theta) = \pi (\mathcal {R} (\theta)) + i \pi (\mathcal {I} (\theta)).
\]

命题 5. --- 设 \(\pi\) 是从 T 到拓扑空间 X 的 \(\mu\) -逆紧映射，\(f\) 是从 X 到拓扑空间 F（不论是否 Hausdorff）的映射。要使 \(f\) 是 \(\pi(\mu)\) -可测的，必要且充分的条件是 \(f \circ \pi\) 是 \(\mu\) -可测的。

我们重新考察 Prop. 4 的证明，先从开头处理的特殊情形开始，沿用相同记号；\(g\) 关于测度 \(\pi(\mu) = \nu\) 可测，当且仅当 \(g_{\mathrm{L}}\) 是 \(\nu'\) -可测的（§1, No.~5, Example）；这又等价于 \(g_{\mathrm{L}} \circ \pi' = (g \circ \pi)_{\mathrm{K}}\) 是 \(\mu_{\mathrm{K}}\) -可测的（Ch. V, §6, No.~2, Prop. 3），最后等价于 \(g \circ \pi\) 是 \(\mu\) -可测的（§1, No.~5, Example）。现在转到一般情形，沿用 Prop. 4 证明中的相同记号；\(f\) 是 \(\nu\) -可测的，当且仅当 \(f\) 是 \(\nu_{\alpha}\) -可测的，对于每个 \(\alpha \in \mathrm{A}\)（§1, No.~7, Prop. 8），因而当且仅当 \(f \circ \pi\) 是 \(\mu_{\alpha}\) -可测的，对于每个 \(\alpha \in \mathbf{A}\)（前述特殊情形），最后当且仅当 \(f \circ \pi\) 是 \(\mu\) -可测的（§1, No.~7, Prop. 8）。

推论。--- 设 X 和 Y 是两个拓扑空间，\(\pi\) 是从 T 到 X 的 \(\mu\) -逆紧映射，\(\pi'\) 是从 X 到 Y 的 \(\pi(\mu)\) -逆紧映射。则映射 \(\pi'' = \pi' \circ \pi\) 是 \(\mu\) -逆紧的，并且 \(\pi''(\mu) = \pi'(\pi(\mu))\)（“测度像的传递性”）。

因为 \(\pi''\) 是 \(\mu\) -可测的（Prop. 5）。置 \(\mu' = \pi(\mu)\)；像负荷 \(\pi'(\mu'^{\bullet}) = \pi'(\pi(\mu^{\bullet}))\) 显然等于 \(\pi''(\mu^{\bullet})\)。由于它局部有界，\(\pi''\) 是 \(\mu\) -逆紧的。于是测度 \(\pi''(\mu)\) 和 \(\pi'(\mu')\) 具有相同的本质上积分，故二者相等。

命题 6. --- 设 \(\pi\) 是从 T 到拓扑空间 X 的 \(\mu\) -逆紧映射，B 是 X 的一个 \(\pi(\mu)\) -可测子集。置 A = \(\overline{\pi}^{-1}(B)\)，并令 \(\pi'\) 为从 A 到 B 的映射，它在 A 上与 \(\pi\) 一致。于是 A 是 \(\mu\) -可测的，\(\pi_A\) 和 \(\pi'\) 是 \(\mu_A\) -逆紧的，并且

\[
\left(\pi (\mu)\right) _ {\mathrm{B}} = \left(\pi_ {\mathrm{A}} (\mu_ {\mathrm{A}})\right) _ {\mathrm{B}} = \pi^ {\prime} (\mu_ {\mathrm{A}}).\tag{2}
\]

根据 Prop. 5，集合 A 是 \(\mu\) -可测的（应用于 \(\varphi_{B}\)）；根据诱导测度的定义（No.~1），映射 \(\pi_{A}\) 显然是 \(\mu_{A}\) -可测的，从而 \(\pi'\) 是可测的。令 f 是 \(\mathcal{F}_{+}(B)\) 中的元素；约定上标为 0 表示在 X 和 T 上延拓为 0，则有

\[
\left(\pi (\mu) _ {\mathrm{B}}\right) ^ {\bullet} (f) = \pi (\mu) ^ {\bullet} \left(f ^ {0}\right) = \mu^ {\bullet} \left(f ^ {0} \circ \pi\right) = \mu^ {\bullet} \left(\left(f \circ \pi^ {\prime}\right) ^ {0}\right) = \mu_ {\mathrm{A}} ^ {\bullet} \left(f \circ \pi^ {\prime}\right),
\]

从而 \(\left(\pi(\mu)_{\mathrm{B}}\right)^{\bullet} = \pi'(\mu_{\mathrm{A}}^{\bullet})\)。由于负荷 \(\left(\pi(\mu)_{\mathrm{B}}\right)^{\bullet}\) 局部有界，\(\pi'(\mu_{\mathrm{A}}^{\bullet})\) 也局部有界，因此 \(\pi'\) 是 \(\mu_{A}\) -逆紧的。测度 \(\pi'(\mu_{\mathrm{A}})\) 和 \(\left(\pi(\mu)\right)_{\mathrm{B}}\) 具有相同的本质上积分，故二者相等。另一个关系可用类似方式得到。

命题 7. --- 设 T 是拓扑空间 X 的子空间，i 是从 T 到 X 的注入。

\begin{enumerate}
\def\labelenumi{\alph{enumi})}
\item
  如果 \(\mu\) 是 T 上的测度且 \(i\) 是 \(\mu\) -逆紧的，则测度 \(i(\mu)\) 集中于 T，并且有 \((i(\mu))_{T} = \mu\)。
\item
  如果 \(\lambda\) 是 \(X\) 上的测度且 \(T\) 是 \(\lambda\) -可测的，则 \(i\) 是 \(\lambda_{T}\) -逆紧的，并且 \(i(\lambda_T) = \varphi_T \cdot \lambda\)。
\item
  置 \(\nu = i(\mu)\)；将关系 \(\nu^{\bullet}(\mathbf{A}) = \mu^{\bullet}(\mathbf{A} \cap \mathbf{T})\) 应用于 \(\mathbf{A} = \mathbf{X} - \mathbf{T}\)，表明 \(\nu\) 集中于 \(\mathbf{T}\)。关系 \(\nu_{\mathrm{T}} = \mu\) 是关系 (2) 的一个特殊情形，取 \(\mathbf{B} = \mathbf{T} = \mathbf{A}\)。
\item
  令 \(f\) 是定义在 \(\mathbf{X}\) 上的正函数；置 \(\mu = \lambda_{\mathrm{T}}\)，有 \(\mu^{\bullet}(f \circ i) = \lambda_{\mathrm{T}}^{\bullet}(f_{\mathrm{T}}) = \lambda^{\bullet}(f\varphi_{\mathrm{T}}) \leqslant \lambda^{\bullet}(f)\)（Prop. 1）；因此 \(i\) 是 \(\mu\) -逆紧的。另一方面，\(\mu^{\bullet}(f \circ i)\)（相应地，\(\lambda^{\bullet}(f\varphi_{\mathrm{T}})\)）是 f 关于 \(i(\mu)\)（相应地，\(\varphi_{T} \cdot \lambda\)）的本质上积分。因此这两个测度相等。
\end{enumerate}

注记 6). --- 设 \(\pi\) 是从 T 到拓扑空间 X 的 \(\mu\) -逆紧映射。关于像测度 \(\nu = \pi(\mu)\) 的积分理论可如下归结为 Ch. V, §6 中处理的理论。令 \((\mathrm{K}_{\alpha})_{\alpha \in \mathrm{A}}\)（相应地，\((\mathrm{L}_{\beta})_{\beta \in \mathrm{B}}\)）是 T（相应地，X）关于 \(\mu\)（相应地，\(\nu\)）的一个分割，并置 \(N = T - \bigcup_{\alpha \in A} K_{\alpha}\)，\(P = X - \bigcup_{\beta \in B} L_{\beta}\)。我们可以假设 \(\pi\) 在每个 \(K_{\alpha}\) 上的限制是连续的（§1, No.~8, Prop. 10）。令 \(T', X'\) 是按照 §1, No.~8 的 Scholium 构造的局部紧空间，令 \(\mu'\) 和 \(\nu'\) 是分别与 \(\mu\) 和 \(\nu\) 关联的这些空间上的测度。由于 \(T'\) 的拓扑是子空间 \(K_{\alpha}\) 的拓扑与 N 上离散拓扑的和，\(\pi\) 是从 \(T'\) 到 X 的连续映射，而关系 \(\mu'^{\bullet}(g \circ \pi) = \mu^{\bullet}(g \circ \pi) = \nu^{\bullet}(g)\)（对于 \(g \in \mathcal{F}_{+}(X)\)）表明 \(\pi\) 是 \(\mu'\) -逆紧的且 \(\pi(\mu') = \nu\)。另一方面，恒等映射 \(i\) 是从 X 到 \(X'\) 的 \(\nu\) -逆紧映射，并且 \(i(\nu) = \nu'\)。因此 \(\pi\) 是一个\(\mu'\) -逆紧的从 \(T'\) 到 \(X'\) 的映射，并且 \(\mu'\) 在 \(\pi\) 下的像是 \(\nu'\)（Cor. of Prop. 5）。把 Ch. V, §6 的结果转写出来的工作留给读者。

\subsection{测度的提升}

命题 8. --- 设 T 和 X 是两个拓扑空间，\(\pi\) 是从 T 到 X 的映射。

\begin{enumerate}
\def\labelenumi{\alph{enumi})}
\item
  设 \(\nu\) 是 X 上的有界测度。要使存在 T 上的测度 \(\mu\)，使得 \(\pi\) 是 \(\mu\) -逆紧的且 \(\pi(\mu) = \nu\)，必要且充分的条件是：对于每个数 \(\varepsilon > 0\)，存在紧集 \(\mathrm{K}_{\varepsilon} \subset \mathrm{T}\)，使得 \(\pi\) 在 \(\mathrm{K}_{\varepsilon}\) 上的限制是连续的，并且 \(\nu^{\bullet}(\mathrm{X} - \pi(\mathrm{K}_{\varepsilon})) < \varepsilon\)。
\item
  假设 \(\pi\) 是单射；令 \(\mu\) 和 \(\mu'\) 是 T 上的两个测度，使得 \(\pi\) 对 \(\mu\) 和 \(\mu'\) 都是逆紧的，并且 \(\pi(\mu) = \pi(\mu')\)。则 \(\mu = \mu'\)。
\end{enumerate}

所述条件 \(a\) 是必要的。事实上，若 \(\pi\) 是 \(\mu\) -逆紧的且 \(\pi(\mu) = \nu\)，则关系 \(\mu^{\bullet}(1) = \nu^{\bullet}(1) < +\infty\) 蕴含 \(\mu\) 有界。将 §1, No.~2 的 Prop. 2 应用于函数 1，得到 T 中存在紧子集 K，使得 \(\mu^{\bullet}(\mathrm{T} - \mathrm{K}) < \varepsilon/2\)。由于 \(\pi\) 是 \(\mu\) -可测的，存在紧集 \(\mathrm{K}_{\varepsilon} \subset \mathrm{K}\)，使得 \(\pi\) 在 \(\mathrm{K}_{\varepsilon}\) 上的限制连续，并且 \(\mu^{\bullet}(\mathrm{K} - \mathrm{K}_{\varepsilon}) < \varepsilon/2\)。于是（No.~3, Prop. 4）

\[
\nu^ {\bullet} \big (\mathrm{X} - \pi (\mathrm{K} _ {\varepsilon}) \big) = \mu^ {\bullet} \big (\mathrm{T} - \pi^{-1} \big (\pi (\mathrm{K} _ {\varepsilon}) \big) \big) <   \varepsilon .
\]

为证明条件是充分的，我们先处理一个特殊情形。

引理 1. --- 设 U 和 V 是两个紧空间，h 是从 U 到 V 的连续满射。于是映射 \(\lambda \mapsto h(\lambda)\) 从 \(\mathcal{M}_{+}(U)\) 到 \(\mathcal{M}_{+}(V)\) 是满射。

事实上，令 \(a\) 为线性映射 \(f \mapsto f \circ h\)，它从 \(\mathcal{C}(\mathrm{V})\) 映到 \(\mathcal{C}(\mathrm{U})\)；由于 \(h\) 是满射，\(a\) 是从 \(\mathcal{C}(\mathrm{V})\) 到子空间 \(\mathrm{H}\)（\(\mathcal{C}(\mathrm{U})\) 的子空间）的等距映射。令 \(\theta\) 是 \(\mathrm{V}\) 上的正测度；于是 \(\theta \circ a^{-1}\) 是 \(\mathrm{H}\) 上的连续线性形式，可延拓为 \(\eta\)（定义在 \(\mathcal{C}(\mathrm{U})\) 上）且具有相同范数。于是 \(\eta\) 是 \(\mathrm{U}\) 上的测度，并且 \(\theta(f) = \eta(f \circ h)\) 对所有 \(f \in \mathcal{C}(\mathrm{V})\) 成立，所以 \(\theta = h(\eta)\)。最后，\(\theta(1) = \| \theta \| = \| \eta \|\)，且 \(\theta(1) = \eta(1)\)，故 \(\eta\) 是正的（Ch. V, §5, No.~5, Prop. 9）。

现在证明条件 \(a\) ）的充分性。该条件蕴含存在 \((\mathbf{K}_n)_{n \geqslant 1}\) 的紧子集序列，定义在 \(T\) 中，使得 \(\pi\) 在每个 \(\mathbf{K}_n\) 上的限制连续，并且对于每个 \(n\) 都有 \(\nu^{\bullet}(\mathrm{X} - \pi(\mathrm{K}_n)) < 1/n\)。可以假设序列 \((\mathbf{K}_n)\) 是递增的。置 \(\mathrm{L}_n = \pi(\mathrm{K}_n)\)，并以 \(\nu_n'\) 表示 \(\varphi_{\mathrm{L}_n - \mathrm{L}_{n-1}} \cdot \nu_{\mathrm{L}_n}\) 上的测度 \(\mathrm{L}_n\)，约定 \(\mathrm{L}_0 = \emptyset\)。

由于限制 \(\pi_{\mathrm{K}_n}\) 连续，存在测度 \(\mu_n'\)，它定义在 \(\mathbf{K}_n\) 上，使得 \(\pi_{\mathrm{K}_n}(\mu_n') = \nu_n'\)（Lemma 1）。令 \(\mu_n\) 为 \(\mu_n'\) 在 \(\mathbf{K}_n\) 到 \(\mathbf{T}\) 的典范注入下的像，并令 \(g\) 是 \(\mathcal{F}_{+}(\mathbf{X})\) 中的元素。依次使用 \(\nu\) 集中于 \(\bigcup_{n} \mathrm{L}_n\) 这一事实、§1, No.~5 的 Prop. 4、§1, No.~2 的 Prop. 2、No.~3 的 Prop. 4 以及最后 No.~3 的 Prop. 7，可得

\[
\begin{array}{r l} \nu^ {\bullet} (g) & = \sum_ {n} \nu^ {\bullet} (\varphi_ {\mathrm{L} _ {n} - \mathrm{L} _ {n - 1}} g) = \sum_ {n} \nu_ {n} ^ {\prime   \bullet} (g _ {\mathrm{L} _ {n}}) = \sum_ {n} \mu_ {n} ^ {\prime   \bullet} (g _ {\mathrm{L} _ {n}} \circ \pi_ {\mathrm{K} _ {n}}) \\ & = \sum_ {n} \mu_ {n} ^ {\prime   \bullet} ((g \circ \pi) _ {\mathrm{K} _ {n}}) = \sum_ {n} \mu_ {n} ^ {\bullet} (g \circ \pi). \end{array}
\]

在这个公式中取 \(g = 1\)，可见族 \((\mu_n)\) 可求和，且其和是一个有界测度 \(\mu\)（§1, No.~7, Prop. 7）。根据 No.~3 的 Prop. 5，对于所有 \(\pi\)，映射 \(\mu_n\) 是 \(n\) -可测的，因为 \(\pi_{\mathrm{K}_n}\) 连续，因而是 \(\mu_n'\) -可测的；由此 \(\pi\) 是 \(\mu\) -可测的（§1, No.~7, Prop. 8），又因 \(\mu\) 有界而是 \(\mu\) -逆紧的。上述关系于是证明测度 \(\pi(\mu)\) 与 \(\nu\) 具有相同的本质上积分，故二者相等（§1, No.~2, Cor. of Prop. 2）。

最后，假设 \(\pi\) 是单射，并证明 \(b\)。令 \(f\) 是 \(\mathcal{F}_{+}(\mathrm{T})\) 中的元素；由于 \(\pi\) 是单射，存在函数 \(g \in \mathcal{F}_{+}(\mathrm{X})\)，使得 \(f = g \circ \pi\)，并且置 \(\nu = \pi(\mu) = \pi(\mu')\) 后，由 No.~3 的 Prop. 4 有

\[
\mu^ {\bullet} (f) = \mu^ {\bullet} (g \circ \pi) = \nu^ {\bullet} (g) = \mu^ {\prime \bullet} (g \circ \pi) = \mu^ {\prime \bullet} (f).
\]

因此测度 \(\mu\) 和 \(\mu'\) 具有相同的本质上积分，这蕴含二者相等（§1, No.~2, Cor. of Prop. 2）。

注记。--- 假设 \(\pi\) 是单射。设 \(\theta\) 是复测度，\(\pi\) 是 \(\theta\) -逆紧的且 \(\pi(\theta)=0\)；则 \(\theta=0\)。事实上，将 \(\theta\) 分解为实部和虚部后，可以归结为 \(\theta\) 为实测度的情形。此时有 \(\pi(\theta^{+})=\pi(\theta^{-})\)，因而 \(\theta^{+}=\theta^{-}\)（Prop. 8），最后 \(\theta=0\)。

下面是 Prop. 8 的条件 \(a)\) 总是满足的一个重要情形。

命题 9. --- 设 T 是苏斯林空间（GT, IX, §6, No.~2, Def. 2），X 是 Hausdorff 空间，\(\pi\) 是从 T 到 X 的连续满射，\(\nu\) 是 X 上的有界测度。则存在 T 上的有界测度 \(\mu\)，使得 \(\pi(\mu) = \nu\)。

这些假设显然蕴含 \(\mathbf{X}\) 是苏斯林空间。

考虑集合函数 \(c: \mathrm{A} \mapsto \nu^{\bullet}(\pi(\mathrm{A}))\)，它定义在 \(\mathfrak{P}(\mathrm{T})\) 上。关系 \(\mathrm{A} \subset \mathrm{B}\) 蕴含 \(c(\mathrm{A}) \leqslant c(\mathrm{B})\)；若 \((\mathrm{A}_n)\) 是 \(\mathrm{T}\) 的子集的递增序列，且 \(\mathrm{A} = \bigcup_{n \in \mathbf{N}} \mathrm{A}_n\)，则 \(c(\mathrm{A}) = \sup_n c(\mathrm{A}_n)\)，这是因为 \(\nu^{\bullet}\) 是负荷。最后，令 \(\mathrm{A} \subset \mathrm{T}\)，令 \(\varepsilon\) 是一个 \(> 0\) 的数；取 \(\mathrm{G}\) 为 \(\mathrm{X}\) 中包含 \(\pi(\mathrm{A})\) 的开子集，使得 \(\nu^{\bullet}(\mathrm{G}) \leqslant \nu^{\bullet}(\pi(\mathrm{A})) + \varepsilon\)（§1, No.~9, Prop. 13）；\(\mathrm{H} = \overline{\pi}^{-1}(\mathrm{G})\) 中的开子集 \(\mathrm{T}\) 包含 \(\mathrm{A}\)，并且 \(c(\mathrm{H}) \leqslant c(\mathrm{A}) + \varepsilon\)。因此函数 \(c\) 是 \(\mathrm{T} (\mathrm{TG},\mathrm{IX},\S 6,\mathrm{No.~10},\mathrm{Def.~9})^{(1)}\) 上的右连续容度，而可容性定理（loc. cit., Th. 6）蕴含等式 \(c(\mathrm{T}) = \sup_{\mathrm{K}}c(\mathrm{K})\)，其中 \(K\) 遍历 \(T\) 的紧子集。于是 Prop. 8 蕴含所需测度 \(\mu\) 的存在性。

\subsection{两个测度的乘积}

设 S 和 T 是两个拓扑空间，分别赋有两个（正的）预测度 \(\lambda\) 和 \(\mu\)，令 X 为乘积空间 \(S \times T\)。设 K 是 X 的紧子集；以 A 和 B 分别表示 K 在 S 和 T 上的投影，并置

\[
\nu_ {\mathrm{K}} = \left(\lambda_ {\mathrm{A}} \otimes \mu_ {\mathrm{B}}\right) _ {\mathrm{K}}.\tag{3}
\]

由此在 X 上定义了一个预测度。事实上，令 L 是 X 中包含 K 的紧子集，令 C 和 D 是它的两个投影；则 \(A \subset C\)、\(B \subset D\)，因而利用诱导测度的传递性以及 Ch. V, §8, No.~5 的 Prop. 12，有

\[
\begin{array}{r l} & {(\nu_ {\mathrm{L}}) _ {\mathrm{K}} = \big ((\lambda_ {\mathrm{C}} \otimes \mu_ {\mathrm{D}}) _ {\mathrm{L}} \big) _ {\mathrm{K}} = (\lambda_ {\mathrm{C}} \otimes \mu_ {\mathrm{D}}) _ {\mathrm{K}}} \\ & {\qquad = \big ((\lambda_ {\mathrm{C}} \otimes \mu_ {\mathrm{D}}) _ {\mathrm{A} \times \mathrm{B}} \big) _ {\mathrm{K}} = (\lambda_ {\mathrm{A}} \otimes \mu_ {\mathrm{B}}) _ {\mathrm{K}} = \nu_ {\mathrm{K}}.} \end{array}
\]

定义 6. --- 由 (3) 定义的预测度 \(\nu\) 称为 \(\lambda\) 和 \(\mu\) 的乘积预测度，记作 \(\lambda \otimes \mu\)。

这个定义显然可以推广到 \(\lambda\) 和 \(\mu\) 是复预测度的情形；此时有 \(|\lambda\otimes\mu|=|\lambda|\otimes|\mu|\)（Ch. III, §4, No.~2, Prop. 3 and Ch. IV, §5, No.~7, Lemma 3）。

若 \(f\) 和 \(g\) 分别是定义在 \(\overline{\mathbf{R}}_+\) 和 \(\mathbf{C}\) 上的函数，则函数 \((s,t)\mapsto f(s)g(t)\) 定义在 \(S\times T\) 上，记作 \(f\otimes g\)。

命题 10. --- 设 \(\nu\) 是 \(\lambda\) 和 \(\mu\) 的乘积预测度；对于每个函数 \(f \in \mathcal{F}_{+}(\mathrm{S})\) 和每个函数 \(g \in \mathcal{F}_{+}(\mathrm{T})\)，

\[
\nu^ {\bullet} (f \otimes g) = \lambda^ {\bullet} (f) \mu^ {\bullet} (g).\tag{4}
\]

预测度 \(\nu\) 是满足 (4) 的 \(S \times T\) 上唯一的预测度。

当 K（相应地，L）遍历 S（相应地，T）的紧子集时，有

\[
\begin{array}{l} \nu^ {\bullet} (f \otimes g) = \sup _ {K, L} \nu_ {K \times L} ^ {\bullet} \big ((f \otimes g) _ {K \times L} \big) = \sup _ {K, L} (\lambda_ {K} \otimes \mu_ {L}) ^ {\bullet} (f _ {K} \otimes g _ {L}) \\ \qquad = \sup _ {K, L} \lambda_ {K} ^ {\bullet} (f _ {K}) \cdot \mu_ {L} ^ {\bullet} (g _ {L}) = \big (\sup _ {K} \lambda_ {K} ^ {\bullet} (f _ {K}) \big) \big (\sup _ {L} \mu_ {L} ^ {\bullet} (g _ {L}) \big) \\ \qquad = \lambda^ {\bullet} (f) \mu^ {\bullet} (g) \end{array}
\]

这是由 Ch. V, §8, No.~3 的 Prop. 8 得到的。

设 \(\eta\) 是 \(S \times T\) 上满足 (4) 的另一个预测度，设 \(K\) 和 \(L\) 分别是 \(S\) 和 \(T\) 的紧子集，\(f\) 和 \(g\) 分别是 \(\mathcal{F}_{+}(K)\) 和 \(\mathcal{F}_{+}(L)\) 中的元素。对于延拓为 0 的函数，有关系 \((f \otimes g)^{0} = f^{0} \otimes g^{0}\)，因此（§1, No.~2, Prop. 2）

\[
\begin{array}{r l} & {\eta_ {\mathrm{K} \times \mathrm{L}} ^ {\bullet} (f \otimes g) = \eta^ {\bullet} \big ((f \otimes g) ^ {0} \big) = \eta^ {\bullet} (f ^ {0} \otimes g ^ {0})} \\ & {\qquad = \lambda^ {\bullet} (f ^ {0}) \mu^ {\bullet} (g ^ {0}) = \lambda_ {\mathrm{K}} ^ {\bullet} (f) \mu_ {\mathrm{L}} ^ {\bullet} (g).} \end{array}
\]

特别地，若取 \(f \in \mathcal{K}_{+}(\mathrm{K})\)、\(g \in \mathcal{K}_{+}(\mathrm{L})\)，则可见 \(\eta_{\mathrm{K} \times \mathrm{L}}\) 具有乘积测度 \(\lambda_{\mathrm{K}} \otimes \mu_{\mathrm{L}}\) 的特征性质（Ch. III, §4, No.~1, Th. 1）。因此 \(\eta_{\mathrm{K} \times \mathrm{L}} = \nu_{\mathrm{K} \times \mathrm{L}}\)；由于 \(S \times T\) 的每个紧子集都包含于某个形如 \(K \times L\) 的集合中，诱导测度的传递性蕴含 \(\eta = \nu\)。

推论 1. --- 如果 \(\lambda\) 和 \(\mu\) 是测度，则 \(\nu\) 是测度。

事实上，令 \(x = (s, t)\) 是 X 中的一点，令 U 和 V 分别是 s 和 t 的邻域，并满足 \(\lambda^{\bullet}(\mathrm{U}) < +\infty\)、\(\mu^{\bullet}(\mathrm{V}) < +\infty\)；集合 \(U \times V\) 是 x 的邻域，并且由 (4) 有 \(\nu^{\bullet}(\mathrm{U} \times \mathrm{V}) = \lambda^{\bullet}(\mathrm{U}) \mu^{\bullet}(\mathrm{V}) < +\infty\)；因此负荷 \(\nu^{\bullet}\) 局部有界，预测度 \(\nu\) 是测度。

这个结果立即推广到复测度。

推论 2. --- 如果 A 是 S 的对 \(\lambda\) 局部可忽略的子集，则 \(A \times T\) 对 \(\nu\) 局部可忽略。

推论 3. --- 假设 \(\lambda\)（相应地，\(\mu\)）是 S（相应地，T）上测度的可求和族 \((\lambda_{\alpha})_{\alpha \in A}\)（相应地，\((\mu_{\beta})_{\beta \in B}\)）之和。则族 \((\lambda_{\alpha} \otimes \mu_{\beta})_{(\alpha, \beta) \in A \times B}\) 可求和，且其和为 \(\lambda \otimes \mu\)。

事实上，令 \(p\) 为负荷 \(\sum_{\alpha, \beta} (\lambda_{\alpha} \otimes \mu_{\beta})^{\bullet}\)；若 \(f \in \mathcal{F}_{+}(S)\) 且 \(g \in \mathcal{F}_{+}(T)\)，则显然有 \(p(f \otimes g) = \lambda^{\bullet}(f)\mu^{\bullet}(g)\)。推论 1 的证明表明 \(p\) 局部有界，因此族 \((\lambda_{\alpha} \otimes \mu_{\beta})\) 可求和（§1, No.~7, Prop. 7）。其和 \(\eta\) 满足 \(\eta^{\bullet} = p\)（§1, No.~7, Prop. 7），而 Prop. 10 蕴含 \(\eta = \nu\)。

\subsection{关于两个测度乘积的积分}

在本号中，\(\lambda\) 和 \(\mu\) 分别表示 S 和 T 上的测度，\(\nu\) 表示乘积测度 \(\lambda \otimes \mu\)，定义在 \(S \times T\) 上。此外，若 \(f\) 是 \(S \times T\) 上的正函数，则对于每个 \(s \in S\)，以 \(f_s\) 表示 T 上的函数 \(t \mapsto f(s,t)\)，以 \(I_f\) 表示 S 上的函数 \(s \mapsto \mu^\bullet(f_s)\)。

引理 2. --- 设 \(f\) 是 \(\nu\) -可测的正函数，定义在 \(S \times T\) 上；对于每个紧子集 \(L\)（属于 \(T\)），令 \(I_f^L\) 为函数 \(s \mapsto \mu^\bullet(f_s \varphi_L)\)，定义在 \(S\) 上。则函数 \(I_f^L\) 是 \(\lambda\) -可测的，并且

\[
\mathrm{I} _ {f} = \sup _ {\mathrm{L}} \mathrm{I} _ {f} ^ {\mathrm{L}}\tag{5}
\]

\[
\nu^ {\bullet} (f) = \sup _ {L} \lambda^ {\bullet} (I _ {f} ^ {L}),\tag{6}
\]

其中 L 遍历 T 的紧子集所成的集合。

首先注意到包含关系 \(L \subset L'\) 蕴含 \(I_{f}^{L} \leqslant I_{f}^{L'}\)；另一方面，\(\mathrm{I}_{f}^{\mathrm{L}}(s) = \mu_{\mathrm{L}}^{\bullet}((f_{s})_{\mathrm{L}})\) 对所有 \(s \in S\) 成立。因此，公式 (5) 立即由 §1, No.~2 中给出的负荷 \(\mu^{\bullet}\) 的定义得到。若 K 是 S 的紧子集，L 是 T 的紧子集，则按构造有 \(\nu_{\mathrm{K} \times \mathrm{L}} = \lambda_{\mathrm{K}} \otimes \mu_{\mathrm{L}}\)，而 Ch. V, §8, No.~3 的 Prop. 7 蕴含关系

\[
\nu^ {\bullet} (f \varphi_ {\mathrm{K} \times \mathrm{L}}) = \lambda_ {\mathrm{K}} ^ {\bullet} \big ((\mathrm{I} _ {f} ^ {\mathrm{L}}) _ {\mathrm{K}} \big).\tag{7}
\]

此外，\(S \times T\) 的每个紧子集都包含于某个形如 \(K \times L\) 的紧集；在前一个公式中对所有 K 和 L 取上包络，得到

\[
\nu^ {\bullet} (f) = \sup _ {L} \sup _ {K} \lambda_ {K} ^ {\bullet} \big ((I _ {f} ^ {L}) _ {K} \big) = \sup _ {L} \lambda^ {\bullet} (I _ {f} ^ {L}),\tag{8}
\]

即得 (6)。

最后，Ch. V, §8, No.~3 的 Prop. 7 蕴含 \(I_{f}^{L}\) 在 S 的每个紧子集 K 上的限制是 \(\lambda_{K}\) -可测的；这等价于说 \(I_{f}^{L}\) 是 \(\lambda\) -可测的。

命题 11. --- 设 \(f\) 是 \(\geqslant 0\) 的下半连续函数，定义在 \(X = S \times T\) 上。

\begin{enumerate}
\def\labelenumi{\alph{enumi})}
\item
  函数 \(f_{s}: t \mapsto f(s, t)\) 在 \(\mathbf{T}\) 上下半连续，对每个 \(s \in S\) 均如此。
\item
  函数 \(I_{f}: s \mapsto \int^{\bullet} f(s, t) d\mu(t)\) 在 S 上下半连续，并且
\end{enumerate}

\[
\iint_ {\mathrm{X}} ^ {\bullet} f (s, t) d \nu (s, t) = \int_ {\mathrm{S}} ^ {\bullet} d \lambda (s) \int_ {\mathrm{T}} ^ {\bullet} f (s, t) d \mu (t).\tag{9}
\]

性质 \(a\) 显然，因为映射 \(t \mapsto f(s,t)\) 从 T 到 \(\overline{\mathbf{R}}\) 是 f 与连续映射 \(f\) \(t \mapsto (s,t)\) 从 T 到 X 的复合。为证明 \(b\)，我们将使用如下引理：

引理 3. --- 设 X 是拓扑空间（不论是否 Hausdorff），f 是定义在 X 上的 \(\geqslant 0\) 下半连续函数；则 f 是 X 上下半连续函数的一个递增序列 \((f_{n})_{n\in\mathbf{N}}\) 的极限，并且其中每个函数 \(f_{n}\) 都是开集特征函数的正系数线性组合。

给定两个整数 \(k \geqslant 1\) 和 \(n \geqslant 1\)，令 \(J_{kn}\) 表示区间 \(]k/2^n, +\infty]\)（在 \(\overline{\mathbf{R}}\) 中）的特征函数。对于每个 \(x \in \overline{\mathbf{R}}_+\)，置 \(u_n(x) = 2^{-n} \sum_{k=1}^{n \cdot 2^n} J_{kn}(x)\)；立即可见，序列 \((u_n(x))_{n \geqslant 1}\) 递增并以 \(x\) 为极限。因此函数序列 \(f_n = u_n \circ f\) 递增并收敛于 \(f\)，并且有 \(f_{n} = 2^{-n}\sum_{k = 1}^{n\cdot 2^{n}}\varphi_{\mathrm{U}(k,n)}\)，

其中 \(\mathrm{U}(k,n)\) 是 X 的开集 \(\overline{f} (]k / 2^n, + \infty ])\)。

现在证明 \(b\)。函数 \(I_f\) 是递增有向函数族 \(I_f^L\) 的上包络，其中 \(L\) 遍历 \(T\) 的紧子集所成集合（引理 2），所以只需证明函数 \(I_f^L\) 是下半连续的；然后对 \(L\) 取上包络，即可由 (6) 得到公式 (9)（§1, No.~6, Prop. 5）。

于是令 \(\mathcal{H}\) 为正下半连续函数 \(f\)（定义在 \(\mathrm{S} \times \mathrm{T}\) 上）所成的集合，且满足：\(\mathrm{I}_f^{\mathrm{L}}\) 对于每个紧子集 \(\mathrm{L}\)（属于 \(\mathrm{T}\)）都是下半连续的。根据 §1, No.~6 的 Prop. 5，\(\mathcal{H}\) 中每个递增有向集的上确界属于 \(\mathcal{H}\)。由引理 3，因此只需证明开集 \(\mathrm{W}\)（它是 \(\mathrm{S} \times \mathrm{T}\) 的开集）的特征函数属于 \(\mathcal{H}\)。此外，根据 \(\mathrm{S} \times \mathrm{T}\) 上乘积拓扑的定义，开集 \(\mathrm{W}\) 是一个递增有向族 \((\mathrm{W}_{\alpha})_{\alpha \in \mathrm{A}}\) 的并，这些开集形如

\[
\mathrm{W} = \bigcup_ {1 \leqslant i \leqslant n} \left(\mathrm{U} _ {i} \times \mathrm{V} _ {i}\right),
\]

其中 \(U_{i}\) 在 S 中是开集，\(V_{i}\) 在 T 中是开集；根据上面的注记，只需证明这种开集的特征函数属于 H。令 \(s \in S\)，令 U 为由包含 s 的开集 \(U_{i}\) 所成族的交（该族可能为空）；立即可见，\(\varphi_{\mathrm{W}}(s,t) \leqslant \varphi_{\mathrm{W}}(s',t)\) 对所有 \(s' \in U\) 和 \(t \in T\) 成立，从而通过积分，\(\mathrm{I}_{\varphi_{\mathrm{W}}}^{\mathrm{L}}(s) \leqslant \mathrm{I}_{\varphi_{\mathrm{W}}}^{\mathrm{L}}(s')\) 对所有 \(s' \in U\) 成立。因此 \(I_{\varphi_{w}}^{L}\) 下半连续，命题得证。

推论 1. --- 设 \(f\) 是定义在 \(\mathbf{X} = \mathbf{S} \times \mathbf{T}\) 上的正数值函数；则

\[
\iint_ {\mathrm{X}} ^ {*} f (s, t) d \nu (s, t) \geqslant \int_ {\mathrm{S}} ^ {*} d \lambda (s) \int_ {\mathrm{T}} ^ {*} f (s, t) d \mu (t).\tag{10}
\]

事实上，令 \(g\) 是 \(\mathbf{X}\) 上满足 \(g \geqslant f\) 的下半连续函数；根据 Prop. 11，

\[
\begin{array}{r l} \iint_ {\mathrm{X}} ^ {*} g (s, t) d \nu (s, t) & = \iint^ {\bullet} g (s, t) d \nu (s, t) = \int^ {\bullet} d \lambda (s) \int^ {\bullet} g (s, t) d \mu (t) \\ & = \int^ {*} d \lambda (s) \int^ {*} g (s, t) d \mu (t) \geqslant \int^ {*} d \lambda (s) \int^ {*} f (s, t) d \mu (t). \end{array}
\]

对 \(g\) 取下包络，即得不等式 (10)。

推论 2. --- 设 f 是定义在 S × T 上的数值函数且是 ν-可忽略的。则函数 \(f_{s}: t \mapsto f(s, t)\) 对于 λ-几乎每个 \(s \in S\) 都是 μ-可忽略的。

命题 12. --- 设 f 是定义在 X = S × T 上的 ν-可测正函数。假设 f 是 ν-适度的（相应地，μ 是适度的）。则：

\begin{enumerate}
\def\labelenumi{\alph{enumi})}
\item
  集合 N 由那些 \(s \in S\) 组成，使函数 \(f_{s}: t \mapsto f(s, t)\) 不是 \(\mu\) -可测；对于 \(\lambda\) 而言它是可忽略的（相应地，局部可忽略的）。
\item
  映射 \(s \mapsto \int^{\bullet} f(s, t) \, d\mu(t)\) 是 \(\lambda\) -可测的，并且
\end{enumerate}

\[
\iint_ {\mathrm{X}} ^ {\bullet} f (s, t) d \nu (s, t) = \int_ {\mathrm{S}} ^ {\bullet} d \lambda (s) \int_ {\mathrm{T}} ^ {\bullet} f (s, t) d \mu (t).\tag{11}
\]

我们先证明 \(b)\) 在 \(f\) 是 \(\nu\) -适度的情形。根据引理 2，当存在紧子集 \(L\)（属于 \(T\)）使得 \(f\) 在 \(S \times L\) 外为零时，结论的这一部分成立；因为此时有 \(I_f = I_f^{L'}\)；对于每个紧子集 \(L'\)（属于 \(T\) 且包含 \(L\)），公式 (11) 化为 (6)。特别地，\(b)\) 对于函数 \(f\)（该函数在 \(S \times T\) 的某个紧子集外为零）已得证。另一方面，\(b)\) 在 \(f\) 是 \(\nu\) -可忽略的情形成立（Prop. 11 的 Cor. 1）。由于每个 \(\nu\) -适度函数都是一个 \(\nu\) -可忽略函数与一列紧支集函数之和（§1, No.~9, Cor. 3 of Prop. 14），断言 \(b)\) 在 \(f\) 是 \(\nu\) -适度的情况下成立。

同样，\(b)\) 在 \(\mu\) 集中于紧子集 \(L\)（属于 \(T\)）时显然成立（引理 2）。假设 \(\mu\) 是适度的；则存在具有紧支集的测度序列 \((\mu_n)_{n \in \mathbb{N}}\)，定义在 \(T\) 上，使得有关关系 \(\mu = \sum_{n} \mu_n\)。\(\nu = \sum_{n} \lambda \otimes \mu_n\) 对每个测度 \(b)\) 成立，并且 \(\nu_n = \lambda \otimes \mu_n\)；于是 \(\nu = \sum_{n} \nu_n\)。

我们来证明 \(a\)；记 N 为使 \(s \in S\) 中的 \(f_s\) 不是 \(\mu\) -可测的那些点所成的集合；对于 T 的每个紧子集 L，类似地记 \(\mathbf{N}_{\mathrm{L}}\) 为使 \(s \in S\) 中的 \(f_s \varphi_{\mathrm{L}}\) 不是 \(\mu\) -可测的那些点所成的集合。若 K 和 L 分别是 S 和 T 中的紧集，则 \(f_{\mathrm{K} \times \mathrm{L}}\) 关于测度 \(\nu_{\mathrm{K} \times \mathrm{L}} = \lambda_{\mathrm{K}} \otimes \mu_{\mathrm{L}}\) 是可测的，而 Ch. V, §8, No.~2 的 Prop. 2 表明集合 \(\mathbf{N}_{\mathrm{L}}\) 对 \(\lambda_{\mathrm{K}}\) 局部可忽略；由于 K 任意，\(\mathbf{N}_{\mathrm{L}}\) 对 \(\lambda\) 局部可忽略。

假设 f 在形如 \(K \times L\) 的紧集外为零；则 \(N = N_{L}\)，且 N 包含于 K 中；因此 N 是 \(\lambda\) -可忽略的。同样，如果 f 是 \(\nu\) -可忽略的，Prop. 11 的 Cor. 2 蕴含 N 是 \(\lambda\) -可忽略的。然后结合前面两种情形，即可如上处理 f 是 \(\nu\) -适度的情形。

假设 \(\mu\) 集中于 T 的紧子集 L；则仍有 \(N = N_{L}\)，因而 N 对 \(\lambda\) 局部可忽略。由于每个适度测度都是具有紧支集的测度序列之和（\(\S1\)，No.~9, Cor. 5 of Prop. 14），这个结果立即推广到 \(\mu\) 是适度的情形（利用 \(\S1\)，No.~7 的 Prop. 8）。

注记。--- 令 \((\mathrm{K}_{\alpha})_{\alpha\in\mathrm{A}}\) 是 S 关于 \(\lambda\) 的一个分割，并令 \(M = S - \bigcup_{\alpha\in A} K_{\alpha}\)；类似地，定义 \((\mathrm{L}_{\beta})_{\beta\in\mathrm{B}}\) 和 N（针对 T 上的测度 \(\mu\)）。我们把 \(S'\) 记作 S 的子空间 \(K_{\alpha}\) 与离散空间 M 的和所成的局部紧空间；空间 \(T'\) 类似定义，并置 \(X' = S' \times T'\)。局部紧空间 \(X'\) 是 X 的紧子空间族 \((\mathrm{K}_{\alpha} \times \mathrm{L}_{\beta})_{(\alpha,\beta)\in A\times B}\) 与子空间 \(P = (M \times T) \cup (S \times N)\) 的和；后者是 X 的局部 \(\nu\) -可忽略子集（注意 P 一般不是离散空间）。我们在 §1, No.~8 的 Scholium 中看到，存在测度 \(\lambda'\)，定义在 \(S'\) 上，使得关于 \(\lambda\) 和 \(\lambda'\)，可测函数、正函数的本质上积分、本质可积函数及其积分都相同。按照所引用的 Scholium，将测度 \(\mu'\) 定义在 \(T'\) 上，并与 \(\mu\) 关联；将测度 \(\nu'\) 定义在 \(X'\) 上，并与 \(\nu\) 关联；立即可见，\(\nu'\bullet(f \otimes g) = \lambda'\bullet(f)\mu'\bullet(g)\) 对 \(f \in \mathcal{F}_{+}(S)\) 和 \(g \in \mathcal{F}_{+}(T)\) 成立；因此根据 No.~5 的 Prop. 10，有 \(\nu' = \lambda' \otimes \mu'\)。由于 \(X'\) 的拓扑比 X 的拓扑精细，每个 \(\nu\) -适度函数都是 \(\nu'\) -适度的。这个过程使 Lebesgue--Fubini 定理（Ch. V, §8, No.~4, Th. 1）无需新的证明即可推广到当前情形。

\subsection{一个关于测度分解的结果}

命题 13. --- 设 X 是拓扑空间，\(\nu\) 是 X 上的适度测度，p 是从 X 到拓扑空间 T 的 \(\nu\) -逆紧映射，且 \(\mu = p(\nu)\)。假设 X 的每个紧子空间都是可度量的。则存在映射 \(t \mapsto \lambda_{t}\)，从 T 到 \(\mathcal{M}_{+}(X)\)，具有以下性质：

\begin{enumerate}
\def\labelenumi{\alph{enumi})}
\item
  对于每个 \(t \in \mathbf{T}\)，测度 \(\lambda_t\) 集中于 \(\overline{p}^{-1}(t)\)；
\item
  对于 X 上每个普遍可测的 \(^{(1)}\) 正函数 f，函数 \(t \mapsto \lambda_{t}^{\bullet}(f)\) 在 T 上普遍可测，并且
\end{enumerate}

\[
\int_ {\mathrm{X}} ^ {\bullet} f (x) d \nu (x) = \int_ {\mathrm{T}} ^ {\bullet} d \mu (t) \int_ {\mathrm{X}} ^ {\bullet} f (x) d \lambda_ {t} (x);\tag{12}
\]

\begin{enumerate}
\def\labelenumi{\alph{enumi})}
\setcounter{enumi}{2}
\tightlist
\item
  对 \(t \in \mathrm{T}\) 中满足 \(\lambda_t(1) \neq 1\) 的点所成的集合局部 \(\mu\) -可忽略。
\end{enumerate}

此外，如果 \(t \mapsto \lambda_t'\) 是从 \(\mathbf{T}\) 到 \(\mathcal{M}_+(\mathbf{X})\) 的映射，满足条件 \(a)\) 和 \(b)\)，则 \(t \in \mathbf{T}\) 中满足 \(\lambda_t \neq \lambda_t'\) 的点所成的集合局部 \(\mu\) -可忽略。我们需要一个辅助结果：

引理 4. --- 设 X 是拓扑空间，\(\nu\) 是 X 上的测度，\(f\) 是从 X 到拓扑空间 F（不论是否 Hausdorff）的 \(\nu\) -可测映射。则存在从 X 到 F 的普遍可测映射 \(f'\)，它与 \(f\) 局部 \(\nu\) -几乎处处相等。

证明与 Ch. V, §3, No.~4 的 Prop. 7 相同，只需考虑 §1, No.~8 的 Prop. 10。

现在转到 Prop. 13 的证明。

\begin{enumerate}
\def\labelenumi{\Alph{enumi})}
\tightlist
\item
  假设 \(\mathbf{X}\) 是紧的可度量空间，且 \(p\) 是连续满射：
\end{enumerate}

于是空间 T 是紧的可度量空间（GT, IX, §2, No.~10）。根据 Ch. VI, §3, No.~1 的 Th. 1，存在映射 \(t \mapsto \eta_t\)，从 T 到 \(\mathcal{M}_+(X)\)，它淡 \(\mu\) -可测并且标量意义本质 \(\mu\) -可积，使得 \(\nu = \int_{\mathrm{T}} \eta_t d\mu(t)\)，并且 \(\eta_t\) 的总质量为 1 且集中于 \(\overline{p}^{-1}(t)\)，对每个 \(t \in \mathrm{T}\) 均如此。令 \((\mathrm{S}_n)_{n \in \mathbf{N}}\) 是 T 关于 \(\mu\) 的一个分割，使得 H 在每个集合 \(\mathrm{S}_n\) 上的限制连续（§1, No.~8, Props. 10 and 11）；以 \(\Lambda: t \mapsto \lambda_t\) 表示从 T 到 \(\mathcal{M}_+(X)\) 的映射：它在 \(S = \bigcup_{n \in \mathbf{N}} S_n\) 上等于 H，在 T - S 上等于 0。显然，\(\nu = \int_{\mathrm{T}}\lambda_t d\mu (t)\)，且 \(\Lambda\) 满足叙述中的条件 \(a)\)。

设 \(\theta\) 是 T 上的测度；映射 \(\Lambda\) 淡 \(\theta\) -可测且标量意义本质 \(\theta\) -可积，因而也是 \(\theta\) -适配的（Ch. V, §3, No.~1, Prop. 2 b)）。令 \(f\) 是 X 上的普遍可测正函数；将 Ch. V, §3, No.~2 的 Prop. 5 应用于 \(\int \lambda_t d\theta(t)\)，可知映射 \(t \mapsto \lambda_t^{\bullet}(f)\) 是 \(\theta\) -可测的，从而由于 \(\theta\) 的任意性而普遍可测。

公式 (12) 由 Ch. V, §3, No.~2 的 Prop. 5 得到。

\begin{enumerate}
\def\labelenumi{\Alph{enumi})}
\setcounter{enumi}{1}
\tightlist
\item
  假设存在 \(\mathbf{X}'\)，它是 \(\mathbf{X}\) 的紧子集，集中着测度 \(\nu\)，且 \(p_{\mathbf{X}'}\) 连续：
\end{enumerate}

置 \(T' = p(X')\)，\(p' = p_{X'}\)；以 \(\nu'\) 表示测度 \(\nu_{X'}\)，并令 \(\mu'\) 为像测度 \(p'(\nu')\)，定义在 \(T'\) 上。由于 \(p'\) 是连续满射且 \(X'\) 是紧的可度量空间，根据 A)，存在映射 \(\Lambda': t' \mapsto \lambda_{t'}'\)，从 \(T'\) 到 \(\mathcal{M}_{+}(X')\)，满足以下条件：

\(a^{\prime})\) 对于每个 \(t^\prime \in \mathrm{T}^\prime\)，测度 \(\lambda_{t'}'\) 集中于 \(\mathrm{X}' \cap \overline{p}^{-1}(t')\)；

\(b')\) 对于每个普遍可测正函数 \(f'\)（定义在 \(X'\) 上），函数 \(t' \mapsto \lambda_{t'}^{\bullet}(f')\) 在 \(T'\) 上普遍可测，并且

\[
\int_ {\mathrm{X} ^ {\prime}} ^ {\bullet} f ^ {\prime} (x ^ {\prime}) d \nu^ {\prime} (x ^ {\prime}) = \int_ {\mathrm{T} ^ {\prime}} ^ {\bullet} d \mu^ {\prime} (t ^ {\prime}) \int_ {\mathrm{X} ^ {\prime}} ^ {\bullet} f ^ {\prime} (x ^ {\prime}) d \lambda_ {t ^ {\prime}} ^ {\prime} (x ^ {\prime}).
\]

令 \(t \in \mathbf{T}\)；如果 \(t\) 属于 \(\mathbf{T}'\)，则以 \(\lambda_t\) 表示 \(\lambda_t'\) 在 \(\mathbf{X}'\) 到 \(\mathbf{X}\) 的典范注入下的像；如果 \(t\) 属于 \(\mathbf{T} - \mathbf{T}'\)，则置 \(\lambda_t = 0\)。读者不难验证，映射 \(t \mapsto \lambda_t\) 满足叙述中的条件 \(a)\) 和 \(b)\)。

\begin{enumerate}
\def\labelenumi{\Alph{enumi})}
\setcounter{enumi}{2}
\tightlist
\item
  一般情形下的存在性：
\end{enumerate}

由于 X 上的测度 \(\nu\) 是适度的，我们可以取 X 的一个覆盖 \((\mathrm{U}_m)_{m\in \mathbf{N}}\)，它由 \(\nu\) -可积开集组成。此外，令 \((\mathrm{X}_n)_{n\in \mathbf{N}}\) 是 X 关于 \(\nu\) 的一个分割，使得 \(p\) 在每个集合 \(\mathrm{X}_n\) 上的限制连续（§1, No.~8, Props. 10 and 11）；以 \(\nu_n\) 表示 X 上的测度 \(\varphi_{\mathrm{X}_n}\cdot \nu\)，以 \(\mu_n\) 表示它在 \(p\) 下的像。根据 B)，对于每个整数 \(n\in \mathbf{N}\)，存在从 T 到 \(t\mapsto \alpha_t^n\) 的映射 \(\mathcal{M}_{+}(\mathrm{X})\)，满足以下条件：

\(a^{\prime \prime})\) 测度 \(\alpha_{t}^{n}\) 集中于 \(\overline{p}^{-1}(t)\)，对所有 \(t\in \mathbf{T}\) 成立。

\(b^{\prime\prime}\) 如果 f 是 X 上的普遍可测正函数，则 T 上的正函数 \(t \mapsto (\alpha_{t}^{n})^{\bullet}(f)\) 是普遍可测的，并且

\[
\int_ {\mathrm{X}} ^ {\bullet} f (x) d \nu_ {n} (x) = \int_ {\mathrm{T}} ^ {\bullet} d \mu_ {n} (t) \int_ {\mathrm{X}} ^ {\bullet} f (x) d \alpha_ {t} ^ {n} (x).\tag{13}
\]

有 \(\nu = \sum_{n\in \mathbf{N}}\nu_n\) 和 \(\mu = \sum_{n\in \mathbf{N}}\mu_n\)；由 No.~2 的 Prop. 3 及上述引理 4 立即可知，存在 T 上的普遍可测正函数序列 \((g_n)_{n\in \mathbf{N}}\)，使得 \(\mu_{n} = g_{n}\cdot \mu\) 对所有 \(n\in \mathbf{N}\) 成立，并且 \(\sum_{n\in \mathbf{N}}g_n = 1\)。对于每个 \(t\in \mathbf{T}\)，以 \(\beta_t^n\) 表示 X 上的测度 \(g_{n}(t)\cdot \alpha_{t}^{n}\)，以 \(q_{t}\) 表示 X 上的负荷 \(\sum_{n\in \mathbf{N}}(\beta_t^n)^{\bullet}\)。令 \(f\) 是 X 上的普遍可测正函数；利用 No.~2 的 Prop. 2 并对 (13) 中的 \(n\) 求和，得到

\[
\int_ {\mathrm{X}} ^ {\bullet} f (x) d \nu (x) = \int_ {\mathrm{T}} ^ {\bullet} q _ {t} (f) d \mu (t);\tag{14}
\]

此外显然，函数 \(t \mapsto q_t(f)\) 在 \(\mathbf{T}\) 上是普遍可测的。

对于每个 \(m \in \mathbf{N}\)，令 \(\mathbf{E}_m\) 为满足 \(t \in \mathbf{T}\) 且 \(q_t(\mathbf{U}_m) = +\infty\) 的点所成的集合；集合 \(\mathbf{E}_m\) 是普遍可测的，因为映射 \(t \mapsto q_t(\mathbf{U}_m)\) 具有此性质；并且 \(\mathbf{E}_m\) 对 \(\mu\) 局部可忽略，这是将公式 (14) 应用于 \(f = \varphi_{\mathbf{U}_m}\) 且 \(\nu^\bullet(\mathbf{U}_m)\) 有限的结果。于是集合 \(\mathbf{E} = \bigcup_{m \in \mathbf{N}} \mathbf{E}_m\) 普遍可测且对 \(\mu\) 局部可忽略。置 \(\lambda_t = 0\) 对于 \(t \in \mathbf{E}\)。此外，令 \(t \in \mathbf{T} - \mathbf{E}\)；负荷 \(q_t\) 由于开集 \(\mathbf{U}_m\) 覆盖 \(\mathbf{X}\) 且 \(q_t(\mathbf{U}_m)\) 有限而局部有界；根据 §1, No.~7 的 Prop. 7，存在测度 \(\lambda_t\)，定义在 \(\mathbf{X}\) 上，使得 \(q_t = \lambda_t^\bullet\) 且 \(\lambda_t = \sum_{n \in \mathbf{N}} \beta_t^n\)。映射 \(t \mapsto \lambda_t\) 满足叙述中的条件 \(a)\) 和 \(b)\) 是立即的。

D) c) 的证明：

设 \(f\) 是 \(X\) 上普遍可测的正有界函数；我们将证明，函数 \(h_f: t \mapsto \lambda_t^\bullet(f)\) 在 \(T\) 上普遍可测，并且是测度 \(\mu_f = p(f \cdot \nu)\) 关于 \(\mu = p(\nu)\) 的一个密度。令 \(K\) 是 \(T\) 的紧子集，并令 \(A = \overline{p}^{-1}(K)\)。对于每个 \(t \in T\)，测度 \(\lambda_t\) 集中于 \(\overline{p}^{-1}(t)\)；若 \(t\) 属于 \(K\)，则 \(\overline{p}^{-1}(t) \subset A\)，从而 \(\lambda_t^\bullet(f\varphi_A) = \lambda_t^\bullet(f)\)；另一方面，若 \(t\) 属于 \(T - K\)，则 \(\overline{p}^{-1}(t) \subset X - A\)，从而 \(\lambda_t^\bullet(f\varphi_A) = 0\)。将公式 (12) 应用于 \(f \cdot \varphi_A\)，\(^{(1)}\) 可得

\[
\mu_ {f} (\mathbf {K}) = \int_ {\mathrm{A}} ^ {\bullet} f (x) d \nu (x) = \int_ {\mathrm{K}} ^ {\bullet} d \mu (t) \int_ {\mathrm{X}} ^ {\bullet} f (x) d \lambda_ {t} (x) = \int_ {\mathrm{K}} ^ {\bullet} h _ {f} (t) d \mu (t),
\]

这就证明了关系 \(\mu_f = h_f \cdot \mu\)。

令 \(f = 1\)，可见函数 \(h_1: t \mapsto \| \lambda_t \|\) 是测度 \(\mu_1 = \mu\) 关于 \(\mu\) 的一个密度，因此在 T 中局部 \(\mu\) -几乎处处等于 1。

E) 唯一性：

令 \(t \mapsto \lambda_t^i\)（其中 \(i = 1,2\)）是从 \(T\) 到 \(\mathcal{M}_+(X)\) 的两个映射，满足叙述中的条件 \(a)\) 和 \(b)\)。如 C) 中一样，取 X 的一个 \(\mu\) -分割 \((X_n)_{n \in \mathbb{N}}\)，使得 \(X\) \(p_{X_n}\) 对每个 \(n \in \mathbb{N}\) 都是连续的，并置 \(N = X - \bigcup_{n \in \mathbb{N}} X_n\)。对于每个整数 \(n \in \mathbb{N}\)，取 X 上正函数的一个可数集合 \(D_n\)，这些函数在 \(X\) \(X_n\) 外为零，且它们在 \(X_n\) 上的限制在赋范空间 \(\mathcal{C}(X_n)\) 中构成稠密集（将 GT, \(X, \S 3\), No.~3 的 Th. 1 应用于可度量紧空间 \(X_n\)）。置 \(D = \bigcup_{n \in \mathbb{N}} D_n\)。

令 \(f \in \mathbf{D}\)；根据 D)，函数 \(t \mapsto (\lambda_t^1)^{\bullet}(f)\) 和 \(t \mapsto (\lambda_t^2)^{\bullet}(f)\) 是测度 \(\mu_f\) 关于 \(\mu\) 的密度，因此存在 T 中局部 \(\mu\) -可忽略集 \(\mathbf{E}_f\)，使得 \((\lambda_t^1)^{\bullet}(f) = (\lambda_t^2)^{\bullet}(f)\) 对于 \(t \in \mathbf{T} - \mathbf{E}_f\) 成立。此外，根据 (12)，令 \(\mathbf{F}_i\) 为 \(t \in \mathbf{T}\) 中满足 \((\lambda_t^i)^{\bullet}(\mathbf{N}) \neq 0\) 的点所成的集合；它是局部 \(\mu\) -可忽略的，对于 \(i = 1, 2\)。由于 D 是可数的，集合 \(G = \left( \bigcup_{f \in D} E_f \right) \cup F_1 \cup F_2\) 局部 \(\mu\) -可忽略；对于 \(t \in T - G\)，有 \((\lambda_t^1)^{\bullet}(\mathbf{N}) = (\lambda_t^2)^{\bullet}(\mathbf{N}) = 0\)，且

\((\lambda_t^1)_{\mathbf{X}_n} = (\lambda_t^2)_{\mathbf{X}_n}\)，从而根据 §1, No.~8 的 Prop. 9，有 \(\lambda_t^1 = \lambda_t^2\)。

证毕。

注记。--- 1) 如果 X 是苏斯林空间，则 X 的每个紧子空间都是苏斯林空间，因而是可度量的（TG, IX, Appendix I, Cor. 2 of Prop. 3），\(^{(2)}\)，并且 X 上的每个测度 \(\nu\) 都是适度的（§1, No.~9, Remark 1）。根据 Prop. 13，X 上的每个测度因此都相对于每个 \(\nu\) -逆紧映射具有一个分解。

\begin{enumerate}
\def\labelenumi{\arabic{enumi})}
\setcounter{enumi}{1}
\tightlist
\item
  沿用 Prop. 13 的记号，令 \(f\) 是正的 \(\nu\) -可测函数。可以如 Ch. V, §3, No.~2, Prop. 5 中那样证明：\(t \in \mathbf{T}\) 中使 \(f\) 不是 \(\lambda_t\) -可测的点所成的集合局部 \(\mu\) -可忽略，\(t \mapsto \lambda_t^\bullet(f)\) 是 \(\mu\) -可测的，并且关系 (12) 仍然成立。
\end{enumerate}

\section{测度与可加集合函数}

在本节中，我们分别用 \(\mathfrak{K}(\mathrm{T})\) 和 \(\mathfrak{B}(\mathrm{T})\) 表示 Hausdorff 拓扑空间 \(\mathrm{T}\) 的紧子集所成的集合和 \(\mathrm{T}\) 的 Borel σ-集环。

\subsection{测度与紧集的可加集合函数}

定理 1. --- 设 T 是拓扑空间，I 是从 \(\mathfrak{K}(\mathrm{T})\) 到 \(\mathbf{R}_{+}\) 的映射。要使存在 T 上的测度 \(\mu\)，使得 \(\mathrm{I}(\mathrm{K}) = \mu^{\bullet}(\mathrm{K})\) 对所有 \(\mathrm{K} \in \mathfrak{K}(\mathrm{T})\) 成立，必要且充分的条件是 I 满足以下条件：

\begin{enumerate}
\def\labelenumi{\arabic{enumi})}
\item
  如果 K 和 L 是 T 的紧子集且 \(K \subset L\)，则 \(I(K) \leqslant I(L)\)（“I 递增”）。
\item
  如果 K 和 L 是 T 的紧子集，则 I(K ∪ L) ≤ I(K) + I(L)。
\item
  如果 K 和 L 是 T 的不交紧子集，则 \(\mathrm{I}(\mathrm{K}\cup\mathrm{L})=\mathrm{I}(\mathrm{K})+\mathrm{I}(\mathrm{L})\)（“I 可加”）。
\item
  对于 T 的每个紧子集的递减有向族 \((\mathrm{K}_{\alpha})_{\alpha \in \mathrm{A}}\)，都有 \(\mathrm{I}\big(\bigcap_{\alpha \in \mathrm{A}}\mathrm{K}_{\alpha}\big) = \inf_{\alpha \in \mathrm{A}}\mathrm{I}(\mathrm{K}_{\alpha})\)。
\item
  对于每个 \(x \in \mathbf{T}\)，存在一个邻域 \(\mathbf{V}\) 的 \(x\)，使得
\end{enumerate}

\[
\sup_{\substack{\mathbf{K}\in \mathfrak{K}(\mathbf{T})\\ \mathbf{K}\subset \mathbf{V}}}\operatorname{I}(\mathbf{K}) <   + \infty
\]

（“I 局部有界”）。

测度 \(\mu\) 于是唯一。

\(\mu\) 的唯一性由 §1, No.~2 的 Cor. of Prop. 2 得出。上述条件是必要的：前三个条件显然，最后一个条件来自 \(\mu^{\bullet}\) 是局部有界负荷这一事实，而条件 4) 由 §1, No.~6 的 Cor. of Prop. 5 得出。

为证明这些条件是充分的，我们先处理 \(\mathbf{T}\) 为紧空间的情形。

引理 1. --- 假设 T 是紧的，并置 \(l = \mathrm{I}(\mathrm{T})\)。对于每个 \(\mathrm{A} \subset \mathrm{T}\)，置

\[
\mathrm{J}(\mathrm{A}) = \sup_{\substack{\mathrm{K}\in \mathfrak{K}(\mathrm{T})\\ \mathrm{K}\subset \mathrm{A}}}\mathrm{I}(\mathrm{K})\tag{1}
\]

并令 \(\Phi\) 为满足 \(A \subset T\) 且 \(J(A) + J(\mathbb{C}A) = l\) 的集合。于是 \(\Phi\) 是包含 \(\mathfrak{K}(T)\) 的集环，且函数 \(J\) 在 \(\Phi\) 上递增并可加。

显然，\(J\) 是一个递增的集合函数，延拓 \(I\)，并且 \(J(A) + J(\mathbb{C}A) \leqslant l\) 对每个 \(A \subset T\) 都成立。

令 K 和 S 是 T 中的两个紧集；我们首先要证明

\[
\mathrm{J} (\mathrm{K} \cap \mathrm{S}) + \mathrm{J} (\mathbb {C} \mathrm{K} \cap \mathrm{S}) = \mathrm{J} (\mathrm{S}).\tag{2}
\]

考察 I 在 \(\mathfrak{K}(\mathrm{S})\) 上的限制和 J 在 \(\mathfrak{P}(\mathrm{S})\) 上的限制，立即可归结为 \(\mathrm{S} = \mathrm{T}\) 的情形。由于 T 是正规空间，K 是其紧邻域的递减有向族的交，而条件 4) 蕴含：对于每个 \(\varepsilon > 0\)，存在 K 的紧邻域 H，使得 \(\mathrm{I}(\mathrm{H}) \leqslant \mathrm{I}(\mathrm{K}) + \varepsilon\)。令 L 为 T - H 的闭包；L 是紧的，\(\mathrm{L} \cap \mathrm{K} = \varnothing\) 且 \(\mathrm{H} \cup \mathrm{L} = \mathrm{T}\)，因此 \(l = \mathrm{I}(\mathrm{H} \cup \mathrm{L}) \leqslant \mathrm{I}(\mathrm{H}) + \mathrm{I}(\mathrm{L}) \leqslant \mathrm{I}(\mathrm{K}) + \mathrm{I}(\mathrm{L}) + \varepsilon\)（条件 2)），从而关系 \(\mathrm{J}(\mathrm{K}) + \mathrm{J}(\mathbb{C}\mathrm{K}) \geqslant \mathrm{I}(\mathrm{K}) + \mathrm{I}(\mathrm{L}) \geqslant l - \varepsilon\)。由于 \(\varepsilon\) 任意，\(\mathrm{J}(\mathrm{K}) + \mathrm{J}(\mathbb{C}\mathrm{K}) = l\)。这证明了公式 (2)，也证明了包含关系 \(\mathfrak{K}(\mathrm{T}) \subset \Phi\)。

现在证明 \(\Phi\) 是一个集环。由于 \(\Phi\) 显然对取补稳定，只需证明：如果 \(A_{1}\) 和 \(A_{2}\) 是 \(\Phi\) 的元素，则 \(A_{1} \cup A_{2} \in \Phi\)，或者说，

\[
\mathrm{J} \left(\mathrm{A} _ {1} \cup \mathrm{A} _ {2}\right) + \mathrm{J} \left(\mathbf {C} \left(\mathrm{A} _ {1} \cup \mathrm{A} _ {2}\right)\right) \geqslant l.\tag{3}
\]

令 \(\varepsilon\) 是一个 \(>0\) 的数，并且对于 \(i = 1,2\)，令 \(\mathbf{K}_i\) 是包含于 \(\mathbf{A}_i\) 的紧集，令 \(\mathbf{L}_i\) 是包含于 \(\mathbf{C}\mathbf{A}_i\) 的紧集，使得

\[
\mathrm{I} \left(\mathrm{K} _ {i}\right) \geqslant \mathrm{J} \left(\mathrm{A} _ {i}\right) - \varepsilon , \quad \mathrm{I} \left(\mathrm{L} _ {i}\right) \geqslant \mathrm{J} \left(\mathbb {C} \mathrm{A} _ {i}\right) - \varepsilon .
\]

置 \(\mathbf{M}_1 = \mathbf{K}_1\cup \mathbf{L}_1\)；关系 \(l = J(\mathbf{M}_1) + \mathbf{J}(\mathbf{C}\mathbf{M}_1)\) 以及

\[
\mathrm{J} \left(\mathrm{M} _ {1}\right) = \mathrm{I} \left(\mathrm{K} _ {1}\right) + \mathrm{I} \left(\mathrm{L} _ {1}\right) \geqslant \mathrm{J} \left(\mathrm{A} _ {1}\right) + \mathrm{J} \left(\mathbf {C A} _ {1}\right) - 2 \varepsilon = l - 2 \varepsilon
\]

蕴含 \(\mathrm{J}(\mathbb{C}\mathrm{M}_1) \leqslant 2\varepsilon\)。于是若 \(\mathrm{S}\) 是 \(\mathrm{T}\) 的紧子集，关系 (2)（应用于 \(\mathrm{K} = \mathrm{M}_1\)）蕴含 \(\mathrm{J}(\mathrm{S}) \leqslant \mathrm{J}(\mathrm{M}_1 \cap \mathrm{S}) + 2\varepsilon\)，从而

\[
\mathrm{J} (\mathrm{S}) \leqslant \mathrm{J} (\mathrm{K} _ {1} \cap \mathrm{S}) + \mathrm{J} (\mathrm{L} _ {1} \cap \mathrm{S}) + 2 \varepsilon .
\]

将取 \(S = K_{2}\) 和 \(S = L_{2}\) 所得的不等式相加，并考虑不等式 \(\mathrm{J}(K_{2}) + \mathrm{J}(L_{2}) \geqslant l - 2\varepsilon\)，以及 \(K_{1} \cap K_{2}\)、\(L_{1} \cap K_{2}\) 和 \(K_{1} \cap L_{2}\) 是包含于 \(A_{1} \cup A_{2}\) 的三个不交紧集这一事实。以 C 表示这三个紧集的并，则

\[
\begin{array}{r l} & {l - 2 \varepsilon \leqslant \mathrm{J} (\mathrm{K} _ {2}) + \mathrm{J} (\mathrm{L} _ {2}) \leqslant \mathrm{J} (\mathrm{C}) + \mathrm{J} (\mathrm{L} _ {1} \cap \mathrm{L} _ {2}) + 4 \varepsilon} \\ & {\qquad \leqslant \mathrm{J} (\mathrm{A} _ {1} \cup \mathrm{A} _ {2}) + \mathrm{J} \big (\mathbb {C} (\mathrm{A} _ {1} \cup \mathrm{A} _ {2}) \big) + 4 \varepsilon ,} \end{array}
\]

由此立即得到所需公式 (3)，因为 \(\varepsilon\) 是任意的。既然这已经建立，前面的不等式蕴含 \(\mathrm{J}(\mathbf{C}) \geqslant \mathrm{J}(\mathbf{A}_1 \cup \mathbf{A}_2) - 6\varepsilon\)；如果 \(\mathbf{A}_1\) 和 \(\mathbf{A}_2\) 不交，则 \(\mathbf{C}\) 是 \(\mathbf{K}_1 \cap \mathbf{L}_2 \subset \mathbf{A}_1\) 与 \(\mathbf{K}_2 \cap \mathbf{L}_1 \subset \mathbf{A}_2\) 的并，由此可推出 \(\mathrm{J}(\mathbf{A}_1 \cup \mathbf{A}_2) \leqslant \mathrm{J}(\mathbf{A}_1) + \mathrm{J}(\mathbf{A}_2)\)。反向不等式显然成立，故 \(\mathbf{J}\) 确实在 \(\Phi\) 上可加，引理得证。

我们完成 T 为紧空间时的定理证明。令 \(\mathcal{E}(\Phi)\) 为 T 上的 \(\Phi\)-阶梯函数所成的向量空间，赋予一致收敛拓扑（Ch. IV, §4, No.~9, Def. 4）；仍以 J 表示与可加函数 J 关联的 \(\mathcal{E}(\Phi)\) 上的正线性形式（loc. cit., Prop. 18）。由于 \(J(T) = l\)，J 连续且范数为 \(l\)。令 \(\mathcal{H}\) 为 \(\mathcal{E}(\Phi)\) 在一致收敛拓扑下的闭包；立即可验证，J 可由连续性延拓为 \(\mathcal{H}\) 上的正线性形式，仍记为 J。由于 \(\mathcal{H}\) 包含 \(\mathcal{C}(T)\)（loc. cit., No.~10, Prop. 19），J 在 \(\mathcal{C}(T)\) 上的限制是正测度 \(\mu\)。还需证明 \(\mu^{\bullet}(K) = I(K)\)，对于 T 的每个紧子集 K。现在有 \(\mu^{\bullet}(K) = \inf_{f \in S_K} \mu^{\bullet}(f)\)，其中 \(S_K\) 表示 \(\mathcal{C}(T)\) 中满足 \(\geqslant \varphi_K\) 的元素所成的集合（§1, No.~6, Prop. 5）。由于 \(J(f) = \mu^{\bullet}(f)\) 对于 \(f \in \mathcal{C}(T)\) 成立，显然只需证明 \(J(K) \geqslant \inf_{f \in S_K} J(f)\)。如引理 1 的证明中那样，令 H 是 K 的紧邻域，使得 \(J(H) \leqslant J(K) + \varepsilon\)，并令 \(f\) 是 T 上取值介于 0 和 1 之间、在 K 上等于 1 且在 H 外为 0 的连续函数（GT, IX, §4, No.~1, Prop. 1）。于是

\[
\mathrm{J} (f) \leqslant \mathrm{J} (\mathrm{H}) \leqslant \mathrm{J} (\mathrm{K}) + \varepsilon ;
\]

\(\varepsilon\) 任意，所需不等式得证，定理在 \(\mathbf{T}\) 为紧空间时也就得证。

现在转到一般情形。对于 T 中的每个紧集 L，令 \(I_{L}\) 为 I 在 \(\mathfrak{K}(L)\) 上的限制。根据刚处理的特殊情形，存在 L 上唯一的测度 \(\mu_{L}\)，使得 \(\mu_{\mathrm{L}}(\mathrm{K}) = \mathrm{I}_{\mathrm{L}}(\mathrm{K})\) 对每个 \(K \subset L\) 的紧集成立。再令 \(L'\) 是包含于 L 的紧集；有 \((\mu_{\mathrm{L}})_{\mathrm{L}'}^{\bullet}(\mathrm{K}) = \mu_{\mathrm{L}}^{\bullet}(\mathrm{K}) = \mu_{\mathrm{L}'}^{\bullet}(\mathrm{K})\) 对每个 \(K \subset L'\) 的紧集成立，因此 \(\mu_{\mathrm{L}'} = (\mu_{\mathrm{L}})_{\mathrm{L}'}\)；映射 \(\mu : L \mapsto \mu_{L}\) 是一个预测度。条件 5) 表明 \(\mu\) 是测度，而关系 \(\mathrm{I}(\mathrm{K}) = \mu^{\bullet}(\mathrm{K})\) 对所有紧集 \(\mathbf{K} \subset \mathbf{T}\) 显然成立。

注记。--- 1) 在定理 1 的叙述中，条件 4) 可以由下述条件（“右连续性”）替代：

4') 对于每个 \(\mathbf{K} \in \mathfrak{K}(\mathbf{T})\) 和每个 \(\varepsilon > 0\)，存在包含 \(\mathbf{U}\) 的开集 \(\mathbf{K}\)，使得对于每个紧集 \(\mathrm{I}(\mathrm{H}) \leqslant \mathrm{I}(\mathrm{K}) + \varepsilon\)，都有 \(\mathrm{H} \subset \mathbf{U}\)。

事实上，若 \(\mu\) 是测度，则函数 \(I: K \mapsto \mu^{\bullet}(K)\) 满足 \(4'\)（§1, No.~9, Prop. 13）。反之，假设 I 满足 1) 和 \(4'\)；我们来证明 I 于是满足 4)。沿用定理 1 叙述中的记号，取一个 \(\varepsilon > 0\) 以及包含紧集 \(K = \bigcap_{\alpha \in A} K_{\alpha}\) 的开集 U，并使得 \(4'\) 成立。

于是存在指标 \(\beta \in A\)，使得 \(K_{\beta} \subset U\)，这蕴含

\[
\inf _ {\alpha \in \mathbf {A}} \mathrm{I} (\mathrm{K} _ {\alpha}) \leqslant \mathrm{I} (\mathrm{K} _ {\beta}) \leqslant \mathrm{I} (\mathrm{K}) + \varepsilon
\]

且 4) 确实得到验证。

\begin{enumerate}
\def\labelenumi{\arabic{enumi})}
\setcounter{enumi}{1}
\tightlist
\item
  条件 2) 和 3) 的集合可以在定理 1 的叙述中由下述条件替代：
\end{enumerate}

如果 K 和 L 是 T 的紧子集，则

\[
\mathrm{I} (\mathrm{K} \cup \mathrm{L}) + \mathrm{I} (\mathrm{K} \cap \mathrm{L}) = \mathrm{I} (\mathrm{K}) + \mathrm{I} (\mathrm{L}).
\]

事实上，这个条件蕴含 2) 和 3)，另一方面

\[
\mu^ {\bullet} (\mathrm{K} \cup \mathrm{L}) = \mu^ {\bullet} (\mathrm{K} \cap \mathrm{L}) = \mu^ {\bullet} (\mathrm{K}) + \mu^ {\bullet} (\mathrm{L})
\]

对于每个测度 \(\mu\)，这是由特征函数之间的关系 \(\varphi_{\mathrm{K}\cup\mathrm{L}} + \varphi_{\mathrm{K}\cap\mathrm{L}} = \varphi_{\mathrm{K}} + \varphi_{\mathrm{L}}\) 得到的。

\subsection{内正则集合函数}

定义 1. --- 设 T 是拓扑空间，\(\mathfrak{B}(T)\) 是 T 的 Borel σ-集环；令 I 是从 \(\mathfrak{B}(T)\) 到 \(\overline{\mathbf{R}}_{+}\) 的映射。

\begin{enumerate}
\def\labelenumi{\alph{enumi})}
\tightlist
\item
  称 I 是可数可加的，如果对于每个序列 \((\mathbf{A}_{n})\)（其元素是 \(\mathfrak{B}(\mathbf{T})\) 中两两不交的），
\end{enumerate}

\[
\mathrm{I} \left(\bigcup_ {n} \mathrm{A} _ {n}\right) = \sum_ {n} \mathrm{I} (\mathrm{A} _ {n}).\tag{4}
\]

\begin{enumerate}
\def\labelenumi{\alph{enumi})}
\setcounter{enumi}{1}
\tightlist
\item
  称 I 是内正则的，如果对于每个集合 \(\mathbf{A} \in \mathfrak{B}(\mathbf{T})\)，
\end{enumerate}

\[
\mathrm{I} (\mathrm{A}) = \sup _ {\mathrm{K}} \mathrm{I} (\mathrm{K}),\tag{5}
\]

其中 K 遍历 A 的紧子集所成的集合。

\begin{enumerate}
\def\labelenumi{\alph{enumi})}
\setcounter{enumi}{2}
\tightlist
\item
  称 I 有界（相应地，局部有界），如果 \(\mathrm{I(T)} < +\infty\)（相应地，如果 T 中每一点 \(x \in \mathbf{T}\) 都有一个开邻域 V，使得 \(\mathrm{I(V)} < +\infty\)）。
\end{enumerate}

注记。--- 1) 条件 a) 显然蕴含 I 是从 \(\mathfrak{B}(\mathbf{T})\)（按包含关系排序）到 \(\overline{\mathbf{R}}_{+}\) 的递增映射。

\begin{enumerate}
\def\labelenumi{\arabic{enumi})}
\setcounter{enumi}{1}
\item
  假设 I 可数可加；令 \((\mathbf{A}_n)_{n\in \mathbf{N}}\) 是 Borel 集的递增序列，并令 \(\mathbf{A} = \bigcup_{n\in \mathbf{N}}\mathbf{A}_n\)。由于集合 \(\mathrm{D}_0 = \mathrm{A}_0\)、\(\mathrm{D}_n = \mathrm{A}_n - \mathrm{A}_{n - 1}\) 两两不交且其并为 A，有 \(\mathrm{I}(\mathrm{A}) = \sum_{n} \mathrm{I}(\mathrm{D}_{n}) = \lim_{n \to \infty} \mathrm{I}(\mathrm{A}_{n})\)。同样，如果 \((\mathrm{B}_{n})\) 是 Borel 集的递减序列且 \(\mathrm{I}(\mathrm{B}_{0}) < +\infty\)，则 \(\mathrm{I}\left(\bigcap_{n} \mathrm{B}_{n}\right) = \lim_{n \to \infty} \mathrm{I}(\mathrm{B}_{n})\)：只需将前面的结论应用于集合 \(A_{n} = B_{0} - B_{n}\)。
\item
  令 \((\mathbf{A}_n)\) 是 T 的 Borel 集的任意序列。如果 I 可数可加，则 \(\mathrm{I}\left(\bigcup_{n} \mathrm{A}_n\right) \leqslant \sum_{n} \mathrm{I}(\mathrm{A}_n)\)。根据前一个注记，只需对有限序列证明这个不等式。立即可归结为两个集合 \(A_{1}\) 和 \(A_{2}\) 的情形；但关系 (4) 蕴含
\end{enumerate}

\[
\mathrm{I} \left(\mathrm{A} _ {1} \cup \mathrm{A} _ {2}\right) = \mathrm{I} \left(\mathrm{A} _ {1}\right) + \mathrm{I} \left(\mathrm{A} _ {2} - \left(\mathrm{A} _ {1} \cap \mathrm{A} _ {2}\right)\right) \leqslant \mathrm{I} \left(\mathrm{A} _ {1}\right) + \mathrm{I} \left(\mathrm{A} _ {2}\right).
\]

所需不等式立即得到。

\begin{enumerate}
\def\labelenumi{\arabic{enumi})}
\setcounter{enumi}{3}
\item
  如果 I 是可数可加且局部有界的函数，则前一个注记立即蕴含，有 \(\mathrm{I}(\mathbf{K}) < +\infty\)，对于每个紧集 \(\mathbf{K} \subset \mathbf{T}\)。
\item
  可以证明，如果 I 可加，即对有限序列满足 (4)，且 I 内正则，则 I 可数可加（Exer. 7）。读者还可验证，下面 Th. 2 的证明只用到了可加性和内正则性。
\end{enumerate}

定理 2. --- 设 T 是拓扑空间，I 是定义在 \(\mathfrak{B}(T)\) 上、取值于 \(\overline{\mathbf{R}}_{+}\) 的函数。要使存在 T 上的测度 \(\mu\)，使得 \(\mu^{\bullet}(\mathbf{A}) = \mathrm{I}(\mathbf{A})\) 对于每个 \(\mathbf{A} \in \mathfrak{B}(T)\) 成立，必要且充分的条件是 I 可数可加、局部有界且内正则。测度 \(\mu\) 于是唯一。

这三个条件是必要的：事实上，映射 \(A \mapsto \mu^{\bullet}(A)\) 定义在 \(\mathfrak{B}(T)\) 上，是可数可加的（\(\S1\) , No.~5, Cor. of Prop. 4），由测度的定义它是局部有界的（\(\S1\) , No.~2, Def. 5），并且根据 \(\S1\) , No.~2 的注记 3 它是内正则的。

下面证明存在性。显然，I 在 \(\mathcal{K}(\mathrm{T})\) 上的限制满足 Th. 1 叙述中的条件 1)、2)、3) 和 5)；我们还要证明它满足 4)。设 K 是 T 的紧子集，是紧集递减有向族 \((\mathrm{K}_{\alpha})_{\alpha \in \mathrm{A}}\) 的交，并令 \(\varepsilon\) \(>0\)；由于 I 局部有界，存在 K 的一个开（因而是 Borel）邻域 V，使得 \(\mathrm{I(V)} < +\infty\)，于是存在指标 \(\alpha\)，使得 \(\mathrm{K}_{\alpha} \subset \mathrm{V}\)；必要时改变记号，可以假设 \(\mathrm{K}_{\alpha} \subset \mathrm{V}\) 对所有 \(\alpha \in \mathrm{A}\) 成立。由 I 的内正则性，存在紧集 \(\mathrm{L} \subset \mathrm{V} - \mathrm{K}\)，使得 \(\mathrm{I(L)} \geqslant \mathrm{I(V} - \mathrm{K)} - \varepsilon\)；由于 L 不与 K 相交，存在指标 \(\alpha\)，使得 \(\mathrm{L} \cap \mathrm{K}_{\alpha} = \varnothing\)，于是有 \(\mathrm{I}(\mathrm{V} - \mathrm{K}_{\alpha}) \geqslant \mathrm{I}(\mathrm{L}) \geqslant \mathrm{I}(\mathrm{V} - \mathrm{K}) - \varepsilon\)。由于 \(\mathbf{K}_{\alpha} \subset \mathbf{V}\)，可得 \(\mathrm{I}(\mathbf{K}_{\alpha}) \leqslant \mathrm{I}(\mathbf{K}) + \varepsilon\)，条件 4) 得到验证。

根据 Th. 1，存在测度 \(\mu\)，使得 \(\mu^{\bullet}(\mathrm{K}) = \mathrm{I}(\mathrm{K})\) 对所有 \(\mathbf{K} \in \mathfrak{K}(\mathbf{T})\) 成立。集合函数 \(\mu^{\bullet}\) 和 I 在 \(\mathfrak{B}(\mathbf{T})\) 上的内正则性于是蕴含 \(\mu^{\bullet}(\mathbf{A}) = \mathrm{I}(\mathbf{A})\) 对所有 \(\mathbf{A} \in \mathfrak{B}(\mathbf{T})\) 成立，存在性得证。\(\mu\) 的唯一性由 Th. 1 的唯一性断言得出。

\subsection{Radon 空间}

定义 2. --- 设 T 是拓扑空间。如果 T 是 Hausdorff 空间，并且定义在 T 的 Borel σ-集环 \(\mathfrak{B}(\mathrm{T})\) 上、取值于 \(\overline{R}_{+}\) 的每个可数可加且有界（相应地，局部有界）函数都是内正则的，则称 T 是 Radon（相应地，强 Radon）空间。

例如，稍后我们将看到（Prop. 3），每个波兰空间都是强 Radon 空间。特别地，每个具有可数基的局部紧空间都是强 Radon 空间。

存在不是强 Radon 空间的 Radon 空间。

命题 1. --- 每个 Lindelöf \(^{(1)}\) Radon 空间都是强 Radon 空间。

设 T 是 Radon 的 Lindelöf 空间，I 是定义在集环 \(\mathfrak{B}(T)\) 上的正的、可数可加且局部有界的集合函数。满足 \(\mathrm{I}(\mathrm{V}) < +\infty\) 的开集 V 覆盖 T，从中可以选出一个可数覆盖 \((\mathrm{V}_{n})_{n \in \mathbb{N}}\)。置 \(G_{n} = V_{0} \cup V_{1} \cup \cdots \cup V_{n}\) 对每个 \(n \in N\) 成立；置 \(H_{0} = G_{0}\)，并置 \(H_{n} = G_{n} - G_{n-1}\) 对 \(n \geqslant 1\) 成立；最后，令 \(I_{n}\) 为集合函数 \(A \mapsto \mathrm{I}(A \cap H_{n})\)，定义在 \(\mathfrak{B}(T)\) 上，它显然是可数可加且有界的。由于集合 \(H_{n}\) 构成 T 的一个分割，有 \(I = \sum_{n} I_{n}\)。因为空间 T 是 Radon 的，所以对于每个 \(n \in N\)，存在 T 上的有界测度 \(\mu_{n}\)，使得 \(\mu_{n}^{\bullet}(\mathrm{A}) = I_{n}(\mathrm{A})\) 对所有 \(A \in \mathfrak{B}(T)\) 成立；因而也有 \(\sum_{n} \mu_{n}^{\bullet}(\mathrm{A}) = I(\mathrm{A})\)。由于 I 局部有界，族 \((\mu_{n})\) 可求和（§1, No.~7, Prop. 7）；若以 \(\mu\) 表示 \(\sum_{n} \mu_{n}\)，则有 \(\mu^{\bullet}(\mathrm{A}) = I(\mathrm{A})\) 对所有 \(A \in \mathfrak{B}(T)\) 成立，由此 I 是内正则的。换言之，T 是强 Radon 的。

回忆，若拓扑空间 T 的子集 A 对于 T 上的每个测度 \(\mu\) 都是 \(\mu\) -可测的，则称 A 是普遍可测的。这等价于说，对于 T 上每个具有紧支集的测度 \(\mu\)，A 都是 \(\mu\) -可测的（\(\S1\) , No.~8, Prop. 9）。

命题 2. --- 设 X 是拓扑空间，T 是 X 的子空间。

\begin{enumerate}
\def\labelenumi{\alph{enumi})}
\tightlist
\item
  假设 \(\mathbf{T}\) 是 Radon 空间。则对于每个定义在 \(\mathfrak{B}(\mathbf{X})\) 上、正的、可数可加且有界的函数 I，都有
\end{enumerate}

\[
\sup_{\substack{K\text{compact}\\ K\subset T}}I(K) = \inf_{\substack{B\in \mathfrak{B}(X)\\ B\supset T}}I(B).\tag{6}
\]

此外，\(\mathbf{T}\) 在 \(\mathbf{X}\) 中是普遍可测的。

\begin{enumerate}
\def\labelenumi{\alph{enumi})}
\setcounter{enumi}{1}
\tightlist
\item
  反之，假设 X 是 Radon 空间且 T 在 X 中普遍可测；则 T 是 Radon 空间。
\end{enumerate}

我们来证明 \(a\)。记 (6) 的右端为 \(\alpha\)；对于每个 \(n \in \mathbf{N}\)，存在集合 \(C_n \in \mathfrak{B}(X)\)，包含 \(T\)，使得 \(I(C_n) \leqslant \alpha + 2^{-n}\)。置 \(C = \bigcap_{n} C_n\)，则 \(T \subset C\)，且 \(I(C) = \alpha\)。若 \(A \in \mathfrak{B}(T)\)，取 Borel 集 \(B\)（属于 \(X\)），使得 \(A = B \cap T\)（GT, I, §4, No.~5, Prop. 5），并置 \(J(A) = I(B \cap C)\)。这个数与 \(B\) 的选取无关，因为若 \(B'\) 是 \(X\) 中满足 \(A = B' \cap T\) 的另一个 Borel 集，则 \(B \cap C\) 和 \(B' \cap C\) 只相差一个 Borel 集 \(M\)；该集包含于 \(C - T\)，且 \(I(M) = 0\) 由 \(C\) 的构造而成立。显然，\(J(K) = I(K)\) 对每个紧集 \(K \subset T\) 成立。令 \((A_n)\) 是 \(T\) 的 Borel 集的两两不交序列，并且对于每个 \(n\)，令 \(B_n\) 是 \(X\) 的 Borel 集，使得 \(B_n \cap T = A_n\)。必要时将 \(B_n\) 替换为 \(B_n - \left( \bigcup_{k < n} B_k \right)\)，可以假设集合 \(B_n\) 两两不交。置 \(A = \bigcup_{n} A_n\) 和 \(B = \bigcup_{n} B_n\)；于是

\[
\mathrm{J} (\mathrm{A}) = \mathrm{I} (\mathrm{B} \cap \mathrm{C}) = \sum_ {n} \mathrm{I} (\mathrm{B} _ {n} \cap \mathrm{C}) = \sum_ {n} \mathrm{J} (\mathrm{A} _ {n});
\]

因此 J 是 \(\mathfrak{B}(\mathrm{T})\) 上的有界可数可加函数。由于根据假设 T 是 Radon 空间，存在 T 上的有界测度 \(\mu\)，使得 \(\mathrm{J}(\mathrm{A}) = \mu^{\bullet}(\mathrm{A})\) 对每个 \(\mathrm{A} \in \mathfrak{B}(\mathrm{T})\) 成立；于是

\[
\alpha = \mathrm{J} (\mathrm{T}) = \mu^ {\bullet} (\mathrm{T}) = \sup _ {\mathrm{K}} \mu^ {\bullet} (\mathrm{K}) = \sup _ {\mathrm{K}} \mathrm{J} (\mathrm{K}),
\]

根据 \(\mu^{\bullet}\) 的定义，公式 (6) 得证。

现在证明 T 是普遍可测的。令 \(\lambda\) 是 X 上的有界测度；可以把前面的论证应用于集合函数 \(I: A \mapsto \lambda^{\bullet}(A)\)，定义在 \(\mathfrak{B}(X)\) 上，从而存在 T 的紧子集序列 \((K_{n})\)，使得（沿用上面的记号）

\[
\sup _ {n} \lambda^ {\bullet} (\mathrm{K} _ {n}) = \mathrm{J} (\mathrm{T}) = \lambda^ {\bullet} (\mathrm{C}).
\]

置 \(K' = \bigcup_{n \in \mathbf{N}} K_n\)；\(K'\) 是 \(X\) 中的 Borel 集，\(K' \subset T \subset C\)，\(\lambda^\bullet(K') = \lambda^\bullet(C)\)，因此这三个集合只相差 \(\lambda\) -可忽略集，从而 \(T\) 是 \(\lambda\) -可测的。这完成了 \(a\) 的证明。

现在转到 \(b\)。假设 \(X\) 是 Radon 空间且 \(T\) 在 \(X\) 中普遍可测。令 \(I\) 是 \(\mathfrak{B}(T)\) 上的正、有界可数可加函数；于是定义在 \(A \mapsto I(A \cap T)\) 上的函数 \(\mathfrak{B}(X)\) 是正的、有界且可数可加的，因此存在 \(\nu\) 上的有界测度 \(X\)，使得 \(I(A \cap T) = \nu^{\bullet}(A)\) 对所有 \(A \in \mathfrak{B}(X)\) 成立。现在 \(T\) 是 \(\nu\) -可测的；前一个关系表明 \(\nu^{\bullet}(K) = 0\)，其中 \(K\) 遍历 \(X\) 中与 \(T\) 不交的紧集，因此 \(\nu\) 集中于 \(T\)。于是对于 \(A\) 的每个 Borel 集 \(X\)，有 \(I(A \cap T) = \nu^{\bullet}(A \cap T) = \mu^{\bullet}(A \cap T)\)，其中 \(\mu\) 是 \(\nu\) 在 \(T\) 上诱导的测度。最后可得，对于每个 \(I(B) = \mu^{\bullet}(B)\)，有 \(B \in \mathfrak{B}(T)\)，而 \(I\) 确实如此。

推论。--- 如果 X 是 Radon 空间，则 X 的每个 Borel 子集 T 都是 Radon 空间。

因为 \(\mathbf{T}\) 在 \(\mathbf{X}\) 中是普遍可测的。

命题 3. --- 每个苏斯林空间（特别地，每个波兰空间或卢津空间）都是强 Radon 空间。

设 T 是苏斯林空间；由于 T 是 Lindelöf 空间（TG, IX, Appendix I, Cor. of Prop. 1），\(^{(2)}\)，只需证明 T 是 Radon 空间（Prop. 1）。令 I 是定义在 \(\mathfrak{B}(T)\) 上的正的、可数可加且有界函数。我们把 I 延拓到 \(\mathfrak{P}(T)\)，规定对于 T 的每个子集 A，置

\[
\operatorname{I}(\mathbf{A}) = \inf_{\substack{\mathbf{B}\in \mathfrak{B}(\mathbf{T})\\ \mathbf{B}\supset \mathbf{A}}}\operatorname{I}(\mathbf{B}).
\]

我们来证明这个延拓是 T 上的容度（GT, IX, §6, No.~9）。显然，关系 \(\mathbf{A} \subset \mathbf{A}'\) 蕴含 \(\mathrm{I}(\mathbf{A}) \leqslant \mathrm{I}(\mathbf{A}')\)。令 \((\mathbf{A}_n)\) 是 T 的子集的递增序列，并令 \(\mathbf{A} = \bigcup_{n} \mathbf{A}_n\)。包含 \(\mathbf{A}_n\) 的 Borel 集所成的集合族对可数交稳定，因此对于每个 \(n\)，存在 Borel 集 \(\mathbf{B}_n\)，使得 \(\mathbf{A}_n \subset \mathbf{B}_n\) 且 \(\mathrm{I}(\mathbf{A}_n) = \mathrm{I}(\mathbf{B}_n)\)（参见 Prop. 2 的证明）。置 \(\mathbf{C}_n = \bigcap_{p \geqslant n} \mathbf{B}_p\)；\(\mathbf{C}_n\) 是 Borel 集，且 \(\mathbf{A}_n \subset \mathbf{C}_n \subset \mathbf{B}_n\)，因此 \(\mathrm{I}(\mathbf{A}_n) = \mathrm{I}(\mathbf{C}_n)\)。另一方面，序列 \((\mathbf{C}_n)\) 是递增的。令 \(\mathbf{C} = \bigcup_{n} \mathbf{C}_n\)；关系 \(\mathbf{A} \subset \mathbf{C}\) 蕴含

\[
\mathrm{I} (\mathrm{A}) \leqslant \mathrm{I} (\mathrm{C}) = \lim _ {n} \mathrm{I} (\mathrm{C} _ {n}) = \lim _ {n} \mathrm{I} (\mathrm{A} _ {n}),
\]

从而立即得到等式 \(\mathrm{I}(\mathrm{A})=\lim_{n}\mathrm{I}(\mathrm{A}_{n})\)。因此，I 是容度。

如果 \((\mathbf{H}_{n})\) 是 T 中闭集的递减序列，则显然 \(\mathrm{I}\left(\bigcap_{n}\mathrm{H}_{n}\right)=\inf_{n}\mathrm{I}(\mathrm{H}_{n})\)。由此每个 T 的苏斯林子集 F 对于 I 都是可容的（TG, IX, §6, No.~10, Prop. 15）。特别地，T 的每个 Borel 集 A 都是可容的（loc. cit., §6, No.~3, Prop. 10）。\(^{(3)}\) 换言之，

\[
\mathrm{I} (\mathrm{A}) = \sup _ {\mathrm{K}} \mathrm{I} (\mathrm{K}),
\]

其中 K 遍历包含于 A 的紧集；我们已经证明 I 是内正则的。

注记。--- 设 X 是卢津空间（特别地，任意波兰空间），f 是从 X 到（Lusin）正则空间 Y 的双射连续映射。已知（TG, IX, §6, No.~7, Prop. 14），映射 \(B \mapsto f^{-1}(B)\) 是从 Y 的 Borel σ-集环到 X 的 Borel σ-集环的双射。X 和 Y 都是卢津空间，因而都是强 Radon 空间（Prop. 3）。于是立即可知，映射 \(\mu \mapsto f(\mu)\) 是从 X 上有界测度所成的集合到 Y 上有界测度所成的集合的双射。\(^{(4)}\)

\section{测度的逆极限}

在本节中，I 表示一个非空集合，它赋有预序关系，记作 \(i \leqslant j\)，并且关于此关系是有向的。回忆（GT, I, §4, No.~4），由 I 索引的拓扑空间逆系统是一个族 \((\mathrm{T}_i, p_{ij})\)，其中 \(\mathrm{T}_i\) 是拓扑空间，且 \(p_{ij}\) 是从 \(\mathrm{T}_j\) 到 \(\mathrm{T}_i\) 的连续映射，当 \(i \leqslant j\) 时；\(p_{ii}\) 是 \(\mathrm{T}_i\) 的恒等映射；并且 \(p_{ik} = p_{ij} \circ p_{jk}\) 当 \(i \leqslant j \leqslant k\) 时。设 T 是拓扑空间；称族 \((p_i)_{i \in I}\) 为 T 上的相容族，其中映射 \(p_i: \mathrm{T} \to \mathrm{T}_i\) 定义在 T 上，且族 \((p_i)_{i \in I}\) 满足 \(p_i = p_{ij} \circ p_j\) 对于 \(i \leqslant j\)。若 T 中不同的 \(x, y\) 存在 \(i \in I\) 使 \(p_i(x) \neq p_i(y)\)，则映射 \(\mathrm{T} = \varprojlim \mathrm{T}_i\) 将它们分离。拓扑空间 \(p_i\) 称为该逆系统的逆极限；其投影 \(\mathrm{T}_i\) 映射到 \((p_i)_{i \in I}\)，并且这一族投影相容。

\subsection{紧空间上的补集与逆极限}

命题 1. --- 设 X 和 Y 是两个拓扑空间，f 是从 X 到 Y 的连续映射。设 \((\mathrm{K}_{\alpha})_{\alpha \in \mathrm{A}}\) 是 X 的紧子集的递减有向族，其交为 K。则 \(f(\mathrm{K}) = \bigcap f(\mathrm{K}_{\alpha})\) 。

事实上，令 \(y\) 是 \(\bigcap_{\alpha \in \mathbf{A}} f(\mathbf{K}_{\alpha})\) 中的一点；对于每个 \(\alpha \in \mathbf{A}\)，集合 \(\mathbf{L}_{\alpha} = \mathbf{K}_{\alpha} \cap \overline{f}^{-1}(y)\) 是紧的且非空。族 \((\mathbf{L}_{\alpha})_{\alpha \in \mathbf{A}}\) 向下有向，因此其交集 \(\mathbf{L}\) 非空。现在，\(\mathbf{L} = \mathbf{K} \cap \overline{f}^{-1}(y)\)，从而 \(y \in f(\mathbf{K})\)。由此我们证明了包含关系 \(f(\mathbf{K}) \supset \bigcap_{\alpha \in \mathbf{A}} f(\mathbf{K}_{\alpha})\)，反向包含显然。

命题 2. --- 给定由 I 索引的拓扑空间逆系统 \((\mathrm{T}_{i}, p_{ij})\)、拓扑空间 T 以及相容且分离的连续映射族 \(p_{i}: T \to T_{i}\)，则：

\begin{enumerate}
\def\labelenumi{\alph{enumi})}
\item
  对于 T 的每个紧子集 K，有 \(\mathbf{K} = \bigcap_{i\in I}\overline{p_i}^{-1}\big(p_i(\mathbf{K})\big)\) 。
\item
  设 K、L 是 T 的两个不交紧子集。存在 \(i \in I\)，使得对于 \(p_{j}(K)\)，\(p_{j}(L)\) 与 \(j \geqslant i\) 不交。
\item
设 \(x\) 是 \(\bigcap_{i \in I} \overline{p_i}^{-1}(p_i(K))\) 中的一点；对于每个 \(i \in I\)，集合 \(K_i\) 由点 \(y\)（属于 \(K\) 且满足 \(p_i(y) = p_i(x)\)）组成，且是 \(K\) 的非空闭子集。若 \(i \leqslant j\)，则 \(K_i \supset K_j\)；由于 \(K\) 紧，集合 \(\bigcap_{i \in I} K_i\) 非空。令 \(y\) 是 \(\bigcap_{i \in I} K_i\) 中的一点；有 \(y \in K\)，且 \(p_i(y) = p_i(x)\) 对所有 \(i \in I\) 成立，从而 \(y = x\)；最后 \(x \in K\)，这证明了包含关系 \(K \supset \bigcap_{i \in I} \overline{p_i}^{-1}(p_i(K))\)；反向包含显然。
\item
  对每个 \(i \in \mathbf{I}\)，置 \(\mathbf{M}_i = \overline{p_i}^{-1}(p_i(\mathbf{K})) \cap \mathbf{L}\)；这是紧空间 \(\mathbf{L}\) 的闭子集，有 \(\mathbf{M}_i \supset \mathbf{M}_j\) 当 \(i \leqslant j\) 时，并且由 \(a\) )，
\end{enumerate}

\[
\bigcap_ {i \in I} M _ {i} = K \cap L = \varnothing .
\]

因此存在指标 \(i\)，使得 \(\mathbf{M}_i = \varnothing\)。对于 \(j \geqslant i\)，有 \(\overline{p_j}^{-1} (p_j(\mathbf{K})) \cap \mathbf{L} = \mathbf{M}_j \subset \mathbf{M}_i = \varnothing\)，从而 \(p_j(\mathbf{K}) \cap p_j(\mathbf{L}) = \varnothing\)。
